\documentclass[output=paper,colorlinks,citecolor=brown]{langscibook}
\ChapterDOI{10.5281/zenodo.19708122}
\author{Marine Vuillermet\orcid{}\affiliation{University of Zürich} and
        Eva Schultze-Berndt\affiliation{The University of Manchester} and
        Martina Faller\orcid{}\affiliation{The University of Manchester}}
        %esb \orcid{0000-0002-8506-0541}
\title{Introduction: Apprehensional constructions in a cross-linguistic perspective}
\abstract{In this introduction, we discuss recurrent issues of interest raised in the contributions to this volume, and in the wider literature, with respect to apprehensional constructions, broadly defined as constructions which convey a possibility that is also undesirable. The domain of apprehensionality is taken here to encompass constructions at various levels, including main clause modal constructions (apprehensives), subordinate clause constructions akin to English \textit{lest}-clauses (precautioning clauses), and `fear' predicates, which are often found to overlap in their formal manifestation with precautioning clauses. Apprehensional case markers (timitive) are included here only insofar as their marking may overlap with that of precautioning subordinate clauses. We discuss the frequently ambiguous syntactic status of apprehensional clauses and link it to their potential syntactic or pragmatic relationship with a second preemptive clause, encoding an action to be taken to avoid the potential undesirable event. We further outline some interesting language-specific constraints attested for constructions in the apprehensional domain, including constraints on the combination with specific grammatical persons, tense/aspect values, and negation. We also review various semantic and pragmatic issues like the status of the undesirability component (conventional or context-dependent, i.e.\ a pragmatic implicature), and the performative nature of apprehensionals, which links to the well-documented overlap of apprehensionals with prohibitives or optatives in a number of languages. Finally, we examine the remarkable diversity of diachronic origins and patterns of polyfunctionality attested for apprehensionals. This and, more generally, their cross-linguistic semantic and grammatical diversity, can be taken as a correlate of their particularly complex semantics and the interactional nature of their use.}

%move the following commands to the "local..." files of the master project when integrating this chapter
%\usepackage{tabularx}
%\usepackage{langsci-optional}
%\usepackage{langsci-gb4e}
%\bibliography{localbibliography}
%\newcommand{\orcid}[1]{}

\IfFileExists{../localcommands.tex}{
   \addbibresource{../localbibliography.bib}
   \usepackage{tabularx,multicol}
\usepackage{url}
\urlstyle{same}

\usepackage{langsci-optional}
\usepackage{langsci-lgr}
\usepackage{langsci-gb4e}

\babelprovide[import]{hebrew}

\usepackage{paralist} %%may not be needed, currently used in questionnaire


%Siena added packages for introduction; not sure they are needed
\usepackage[normalem]{ulem}
\useunder{\uline}{\ul}{}

%from Treis:
\usepackage{listings}
\lstset{basicstyle=\ttfamily,tabsize=2,breaklines=true}

%%from Francez
% % % \usepackage{cjhebrew}
%\usepackage{tipa}  Also included by D&A, need to check with LSP whether to change it to language science tipa
\usepackage{leipzig}
%%from Dabkowski&AnderBois:

\usepackage[shortlabels]{enumitem} 
\usepackage{centernot}
\usepackage{arydshln}
\usepackage{newunicodechar}
\usepackage{csquotes}
%\let\sups\relax \usepackage{tipa} commented out on 7/3/25 Martina
% \usepackage[all]{nowidow}
\usepackage{listings} \lstset{basicstyle=\ttfamily}
\usepackage{multirow}
\usepackage{rotating}
\usepackage{pdflscape}
\usepackage{xspace}
%\usepackage{tabularx}
%\usepackage{multicol}
\usepackage{supertabular}
\usepackage{hhline}
\usepackage{bbding} %this package is for the checkmarks in Table 6 in Browne et al.

%\usepackage{qtree}
%\usepackage{tree-dvips}

\usepackage{array}
\usetikzlibrary{calc,tikzmark,matrix}
\usepgflibrary{shadings}
\usepackage{pifont}
%from Fedotov:

\usepackage{langsci-textipa}
%\usepackage{multirow}


%\usepackage{xeCJK}

%%%%%%%%%%%%%%%%%%%%%%%%%%%%%%%%%%%%%%%%%%%%%%%%%%%%
%%% EXAMPLES %%%%%%%%%%%%%%%%%%%%%%%%%%%%%%%%%%%%%%%
%%%%%%%%%%%%%%%%%%%%%%%%%%%%%%%%%%%%%%%%%%%%%%%%%%%% 

% if you want the source line of examples to be in
% italics, uncomment the following line
\renewcommand{\exfont}{\itshape}
% to add additional information to the right of
% examples, uncomment the following line

%\usepackage{langsci/styles/jambox}


%%% FOOTNOTE EXAMPLE NUMBERING
% \makeatletter
% \patchcmd{\@footnotetext}{\setcounter{fnx}{0}}{\renewcommand{\thexnumi}{\roman{xnumi}}}{}{}
% \apptocmd{\@footnotetext}{
%     \@noftnotetrue
%     \renewcommand{\thexnumi}{\arabic{xnumi}}%
% }{}{}
%
% \makeatother
\usepackage[linguistics,edges]{forest}
\usepackage{hyphenat}
\usepackage{langsci-branding}

   \makeatletter
\let\thetitle\@title
\let\theauthor\@author
\makeatother

\newcommand{\togglepaper}[1][0]{
   \bibliography{../localbibliography}
   \papernote{\scriptsize\normalfont
     \theauthor.
     \titleTemp.
    To appear in:
    Martina Faller,  Marine Vuillermet \&  Eva Schultze-Berndt (eds.). Apprehensional Constructions in a cross-linguistic perspective. Berlin: Language Science Press. [preliminary page numbering]
   }
   \pagenumbering{roman}
   \setcounter{chapter}{#1}
   \addtocounter{chapter}{-1}
}

\newbool{bookcompile}
\booltrue{bookcompile}
\newcommand{\bookorchapter}[2]{\ifbool{bookcompile}{#1}{#2}}


\newfontfamily\FedotovBoldEmptySetFont{LibertinusSerif-Italic.otf}[AutoFakeBold=2.5]
\newcommand{\FedotovBoldEmpty}{{\FedotovBoldEmptySetFont\bfseries ∅}}

%From Dabkowski&AnderBois


\let\sc\scshape
\let\rm\rmfamily
%\let\it\itshape
\let\tt\ttfamily
\let\nr\normalfont

\newcommand{\wsc}[1]{\textls[80]{\scshape#1}}

\let\textai\textit
% \let\textlc\spacedlowsmallcaps
% \let\textac\spacedallcaps
\let\textlc\textsc
\let\textac\textsc


%%%%%%%%%%%%%%%%%%%%%%%%%%%%%%%%%%%%%%%%%%%%%%%%%%%%
%%% FLAG %%%%%%%%%%%%%%%%%%%%%%%%%%%%%%%%%%%%%%%%%%%
%%%%%%%%%%%%%%%%%%%%%%%%%%%%%%%%%%%%%%%%%%%%%%%%%%%%

% \newcommand{\scott}  [1]{\textcolor{violet}{#1}}
% \newcommand{\maks}   [1]{\textcolor{blue}{#1}}
% % \newcommand{\data}   [1]{\textcolor{teal}{#1}}
% \newcommand{\new}   [1]{#1}
% \newcommand{\newnew}   [1]{\textcolor{teal}{#1}}
% \newcommand{\rev}   [1]{\textcolor{magenta}{#1}}



%%%%%%%%%%%%%%%%%%%%%%%%%%%%%%%%%%%%%%%%%%%%%%%%%%%%
%%% TABLES %%%%%%%%%%%%%%%%%%%%%%%%%%%%%%%%%%%%%%%%%
%%%%%%%%%%%%%%%%%%%%%%%%%%%%%%%%%%%%%%%%%%%%%%%%%%%%

%\newcolumntype{Y}{>{\arraybackslash}p{25.3543464286pt}} % length of L

% \newcolumntype{K}{>{\arraybackslash}p{24.25pt}}
% \newcolumntype{V}{>{\arraybackslash}p{(\textwidth/9-24pt)/2}}
%
\let\ch\relax \newcommand{\ch}[2]{\multicolumn{#1}{c}{\sc#2}}
\let\sx\relax \newcommand{\sx}[2]{\multicolumn{#1}{l}{\emph{-#2}}} % suffix
\let\su\relax \newcommand{\su}{\sx{1}} % suffix unicolumnal
\let\sd\relax \newcommand{\sd}{\sx{2}} % suffix dicolumnal
\let\cx\relax \newcommand{\cx}[2]{\multicolumn{#1}{l}{\emph{=#2}}} % clitic
\let\cu\relax \newcommand{\cu}{\cx{1}} % clitic unicolumnal
\let\cd\relax \newcommand{\cd}{\cx{2}} % suffix dicolumnal
\let\gx\relax \newcommand{\gx}[2]{\multicolumn{#1}{l}{\sc #2}} % gloss
\let\gd\relax \newcommand{\gd}{\gx{2}} % gloss dicolumnal
\let\gr\relax \newcommand{\gr}[1]{\multicolumn{2}{c}{\multirow{2}{*}{\textcolor{gray}{\sc#1}}}}
\let\db\relax \newcommand{\db}[1]{\multicolumn{2}{c}{\multirow{2}{*}{\sc#1}}}



%%%%%%%%%%%%%%%%%%%%%%%%%%%%%%%%%%%%%%%%%%%%%%%%%%%%
%%% EXAMPLES %%%%%%%%%%%%%%%%%%%%%%%%%%%%%%%%%%%%%%%
%%%%%%%%%%%%%%%%%%%%%%%%%%%%%%%%%%%%%%%%%%%%%%%%%%%%

%%% AFFIX BOUNDARY %%%%%%
% \DeclareRobustCommand{\longmn}{\textminus}
% \DeclareRobustCommand{\mn}{-}

%%% CLITIC BOUNDARY %%%%%%
% \DeclareRobustCommand{\longeq}{=}
% \makeatletter\DeclareRobustCommand{\eq}{%
%     \settowidth{\@tempdima}{-}% width of a hyphen
%     \resizebox{\@tempdima}{\height}{=}%
% }\makeatother

%%% FUZZY CLITIC %%%%%%
% \DeclareRobustCommand{\longfc}{%
%     $\overset{\text{\smash{}}}{\text{$\approx$}}$%
% }
% \makeatletter\DeclareRobustCommand{\fc}{%
%     \settowidth{\@tempdima}{-}% width of a hyphen
%     \resizebox{\@tempdima}{\height}{%
%         $\overset{\text{\smash{}}}{\text{$\approx$}}$%
%     }%
% }\makeatother

%%% REDUPLICATION %%%%%%%
% \DeclareRobustCommand{\longtl}{$\sim$}
% \makeatletter\DeclareRobustCommand{\tl}{%
%     \settowidth{\@tempdima}{-}% width of a hyphen
%     \resizebox{\@tempdima}{\height}{$\sim$}%
% }\makeatother

% \newcommand{\mn}          {\textminus}
\newcommand{\src}    [1]  {\jambox[0pt]{\rm(#1)}}
% \newcommand{\ai}     [2]  {\emph{#1} `\textsc{#2}#3'\xspace}
\newcommand{\ai}     [2]  {\emph{#1} `\textsc{#2}'\xspace}
\newcommand{\sane}   [1][]{\ai{=sa'ne}{appr}{#1}}
\newcommand{\srcn}   [1][]{\ai{kûintsû}{srcn}{#1}}
% \newcommand{\mbekane}[1][]{\ai{\eq mbe kañe}{neg.adv aux.inf}{#1}}

% \let\sc\relax \newcommand{\sc}{\normalfont\scshape}
% \let\rm\relax \newcommand{\rm}{\normalfont\rmfamily}
% \let\it\relax \newcommand{\it}{\normalfont\itfamily}

\newcommand{\lb}{{\rm[}}
\newcommand{\rb}{{\rm]}}

\newcommand{\bad}{\makebox[0pt][r]{*\,}}
\newcommand{\infel}{\makebox[0pt][r]{\#\,}}


% \newcommand{\lsqz}[1]{#1}
% \newcommand{\lsqz}[1]{\makebox[0pt][l]{#1}}

% \newcommand{\lsqu}[1]{\makebox[0pt][l]{#1}}
% \newcommand{\rsqu}[1]{\makebox[0pt][r]{#1}}
% \newcommand{\astr}{\makebox[0pt][r]{{\nr*}}}
% \newcommand{\hash}{\makebox[0pt][r]{{\nr\#}}}
% \newcommand{\qstn}{\makebox[0pt][r]{{\nr?}}}
% \newcommand{\bad}{\makebox[0pt][r]{*}}
% \newcommand{\unhappy}{\makebox[0pt][r]{\#}}


%%%%%%%%%%%%%%%%%%%%%%%%%%%%%%%%%%%%%%%%%%%%%%%%%%%%
%%% UNICODE CHARACTERS  %%%%%%%%%%%%%%%%%%%%%%%%%%%%
%%%%%%%%%%%%%%%%%%%%%%%%%%%%%%%%%%%%%%%%%%%%%%%%%%%%

\newunicodechar{⟨}{⟨}
\newunicodechar{⟩}{⟩}

\newcommand{\infix}{⟨}
\newcommand{\outfix}{⟩}



%%%%%%%%%%%%%%%%%%%%%%%%%%%%%%%%%%%%%%%%%%%%%%%%%%%%
%%% IPA %%%%%%%%%%%%%%%%%%%%%%%%%%%%%%%%%%%%%%%%%%%%
%%%%%%%%%%%%%%%%%%%%%%%%%%%%%%%%%%%%%%%%%%%%%%%%%%%%

% \newcommand{\ipa}{\textipa}

\newcommand{\sub}{\textsubscript}
% \let\super\relax \newcommand{\super}{\textsuperscript}

\newcommand{\ts}{\texttslig}
% \let\dz\relax \newcommand{\dz}{\textdzlig}
\newcommand{\tS}{\textteshlig}
\newcommand{\dZ}{\textdyoghlig}
\newcommand{\cj}{\textctc}
\newcommand{\sj}{\textctc}
\newcommand{\zj}{\textctz}
\newcommand{\tcj}{\texttctclig}
\newcommand{\tsj}{\texttctclig}
\newcommand{\dzj}{\textdctzlig}
\newcommand{\gj}{\textbardotlessj}
% \let\nj\relax \newcommand{\nj}{\textltailn}
% \newcommand{\W}{\textturnmrleg}
\newcommand{\wl}{\textturnmrleg}
% \newcommand{\1}{\ipabar{\~\i}{.5ex}{1.1}{}{}}
\newcommand{\dotlessibar}{\ipabar{\i}{.5ex}{1.1}{}{}}
\newcommand{\darkl}{\textltilde}
% \newcommand{\n}{\textsubarch}
% \newcommand{\x}{\textovercross}
\newcommand{\lowered}{\textlowering}
\newcommand{\raised}{\textraising}
\newcommand{\retracted}{\textsubbar}
\newcommand{\unreleased}{\textcorner}
% \newcommand{\syllabic}{\s}
\newcommand{\implosivet}{\texthtt}
\newcommand{\implosived}{\texthtd}
\newcommand{\implosivep}{\texthtp}
\newcommand{\implosiveb}{\texthtb}
\newcommand{\implosivek}{\texthtk}
\newcommand{\implosiveg}{\texthtg}



%%%%%%%%%%%%%%%%%%%%%%%%%%%%%%%%%%%%%%%%%%%%%%%%%%%%
%%% OTHER %%%%%%%%%%%%%%%%%%%%%%%%%%%%%%
%%%%%%%%%%%%%%%%%%%%%%%%%%%%%%%%%%%%%%%%%%%%%%%%%%%%

\newcommand{\ie}{i.\,e.\xspace}
\newcommand{\Ie}{I.\,e.\xspace}
\newcommand{\eg}{e.\,g.\xspace}
\newcommand{\Eg}{E.\,g.\xspace} 




\newcommand{\francoissource}[1]{\hfill\textnormal{{\footnotesize #1} }}
\newcommand{\E}{Ɛ\kern-1.5pt\raisebox{1pt}{̏}\kern1.5pt}
\newcommand{\ù}{{ṵ}\kern-2pt{̀}\kern2pt}
\newcommand{\ȕ}{{ṵ}\kern-2pt{̏}\kern2pt}

\newcommand{\OO}{ɔ\kern-1pt\raisebox{-1pt}{̰}\kern1pt}
\newcommand{\Ǒ}{{{ɔ̌\kern-1pt\raisebox{-.5pt}{̰}\kern1pt}}}


\newcommand{\à}{{a̰}\kern-2pt{̀}\kern2pt}
\newcommand{\ilit}[1]{#1\il{#1}}

   %% hyphenation points for line breaks
%% Normally, automatic hyphenation in LaTeX is very good
%% If a word is mis-hyphenated, add it to this file
%%
%% add information to TeX file before \begin{document} with:
%% %% hyphenation points for line breaks
%% Normally, automatic hyphenation in LaTeX is very good
%% If a word is mis-hyphenated, add it to this file
%%
%% add information to TeX file before \begin{document} with:
%% %% hyphenation points for line breaks
%% Normally, automatic hyphenation in LaTeX is very good
%% If a word is mis-hyphenated, add it to this file
%%
%% add information to TeX file before \begin{document} with:
%% \include{localhyphenation}

\hyphenation{
par-a-digm
com-ple-men-ta-tion
anaph-o-ra
Dor-drecht
Cu-shi-tic
be-ne-dic-tive
pa-ra-digms
Ethi-o-pi-an
hoch-land-ost-ku-schi-ti-schen 
Duu-raa-me
Pa-ken-dorf
Lich-ten-berk
affri-ca-te
affri-ca-tes
com-ple-ments
Eng-lish
}




\hyphenation{
par-a-digm
com-ple-men-ta-tion
anaph-o-ra
Dor-drecht
Cu-shi-tic
be-ne-dic-tive
pa-ra-digms
Ethi-o-pi-an
hoch-land-ost-ku-schi-ti-schen 
Duu-raa-me
Pa-ken-dorf
Lich-ten-berk
affri-ca-te
affri-ca-tes
com-ple-ments
Eng-lish
}




\hyphenation{
par-a-digm
com-ple-men-ta-tion
anaph-o-ra
Dor-drecht
Cu-shi-tic
be-ne-dic-tive
pa-ra-digms
Ethi-o-pi-an
hoch-land-ost-ku-schi-ti-schen 
Duu-raa-me
Pa-ken-dorf
Lich-ten-berk
affri-ca-te
affri-ca-tes
com-ple-ments
Eng-lish
}



   \boolfalse{bookcompile}
   \togglepaper[23]%%chapternumber
}{}



\begin{document}
\maketitle



\section{Introduction}
\label{introintro}

\subsection{Apprehensionality}
\label{subsec:appr}

The functional domain of \textsc{apprehensionality}, broadly speaking, encompasses grammatical markers and constructions which conventionally encode, or pragmatically implicate, that the situation described by the clause is an undesirable possibility. As such, apprehensionality is a subtype of modality. We use the cover term \textit{apprehensional} for a variety of such constructions, illustrated with examples from Northern Australian \ili{Kriol}, \ili{Kalaallisut} and \ili{English} below.\footnote{\citet[1333]{croft2022} uses the term ``apprehensional'' to refer to a semantic relation between two clauses, which actually represents only a subdomain of the functional domain of apprehensionality as defined above. The subdomains are detailed in \sectref{subsec:apprsubtypes}.}


\ea 
\label{ex:intro1} 
\langinfo{Kriol}{\ili{English}-lexified creole; Northern Australia}{\citealt{hodgson2017}}\\
\gll Flowa	im	wetwan.	\textbf{Bambai}	yu	slip	en	boldan. \\
floor is wet \textsc{\textbf{appr}} \textsc{2sg} slip and fall \\
\glt `(Oh no baby girl, don't run around inside!) The floor is wet. You might slip and fall.'
\z

\ea 
\label{ex:intro2} 
\langinfo{Kalaallisut}{Eskimo-Aleut; Greenland}{\citealt[27]{Fortescue1984}}\textup{\footnote{Cited in \citet[31]{Dobrushina2006}.}}\\
\gll Nakka-\textbf{qina}-vutit. \\
fall-\textsc{\textbf{appr}-2sg.ind} \\
\glt `[Be careful], do not fall!'
\z

\ea\label{ex:intro3} He tied his daughter's laces, \textbf{lest} she fall.
\z

Apprehensionals have been labelled ``apprehensive'', ``(ad)monitive'', ``preventative'', ``evitative'', ``timitive'', ``precautioning'' or ``\textit{lest}-markers'', among other terms. Apprehensional clauses are often used to alert or warn the addressee of a potentially undesirable situation as in \REF{ex:intro1} and \REF{ex:intro2}. They may be accompanied by a clause describing a way to avoid the undesirable situation, like `don't run inside' in the context preceding \REF{ex:intro1} or `he tied his daughter's laces' in \REF{ex:intro3}. Following \citet[265]{evans1995}, we label such clauses ``preemptive clauses''.
 
Although wide-spread cross-linguistically, apprehensional markers have received little attention in the linguistics literature. Until recently, their discussion has been mostly restricted to brief sections in reference grammars and other language-specific discussions of some lesser studied languages (e.g.\ in languages from the Australian, Papuan, South Pacific, and Amazonian areas). There exist very few works that address the phenomenon from a cross-linguistic perspective. With the exception of the seminal paper by \citet{Dobrushina2006}, these works tend to focus on an individual language, a small set of languages or a language family (\cite{lichtenberk1995,pakendorfschalley2007,vuillermet2018a}), or a specific subtype of apprehensional (as in the discussion of subordinate `lest' clauses in direct comparison with purpose clauses in \cite[129--144]{schmidtkebode2009}). In contrast to grammatical markers of desirability and volitionality, which are discussed alongside general possibility and necessity modals, markers of undesirability only receive passing mention -- if at all -- in works on modality from a typological-functional perspective (\cite[211]{bybeeetal1994}; \cite[22]{palmer2001}) and have mostly been ignored entirely in the formal semantic literature on modality until very recently \citep{anderbois2020salt,phillips2021,tahar2021, dabkowskianderbois2023}. Even more strikingly, the very existence of apprehensionals has only very recently been acknowledged for colloquial varieties of well-studied languages such as \ili{Dutch} \citep{boogaart2009,boogaart2020} and \ili{German} (\cite{angeloschultzeberndt2018,fallerschultzeberndt2018,fenkoczlon2018}). The description, translation and teaching of all languages as well as theories of modality of all traditions therefore stand to benefit from cross-linguistically informed explorations of this category.

The present volume aims to showcase the variety of linguistic manifestations of apprehensionality, and explore the parameters by which this broader domain can be subcategorised. It includes 12 selected papers from two workshops on the topic of apprehensionality, one at the 2018 European Linguistics Society meeting,\footnote{\url{https://societaslinguistica.eu/wp-content/uploads/2020/03/SLE-2018-Book-of-abstracts.pdf}} co-organised by Schultze-Berndt and Faller, and the other at the 2018 Syntax of the World Languages conference,\footnote{\url{https://swl8.sciencesconf.org/resource/page/id/6.html}} co-organised by Vuillermet and Schultze-Berndt. Each of the papers responds to a set of issues that were raised during the workshops, and presents in-depth discussions of apprehensionals in an individual language, language group or area, by an expert on these languages and based on fieldwork or corpus data. The languages discussed in this volume represent all major linguistic areas and the following language families/subgroups: \ili{Dene}, \ili{A'ingae}, \ili{Mande}, \ili{Semitic}, \ili{East Caucasian} (\ili{Nakh-Daghestanian}), \ili{Slavic}, \ili{Sino-Tibetan} (with \ili{Kiranti} and \ili{Sinitic} subgroups), \ili{Korean}, \ili{Oceanic}, and \ili{Ngumpin-Yapa}. In addition, we have included a questionnaire designed by Vuillermet for the in-depth exploration of variation within apprehensional constructions (based on \cite{VuillermetM2018questionnaire}). The volume is aimed at typologists, semanticists interested in modality, and general linguists, especially those engaged in fieldwork-based research on lesser studied languages.

Some of the main issues addressed by the papers in this volume will be introduced below: \sectref{subsec:apprsubtypes} is a brief and general presentation of the apprehensional subtypes, while \sectref{sec:morph} more specifically targets morphosyntax, \sectref{sec:semantics} semantics, and \sectref{sec:diachrony} diachrony. \sectref{intro:syntacticstatus} expressly explores the complexity of determining the syntactic status of an apprehensional clause, and, consequently, the exact function(s) encoded by the apprehensional morpheme or construction. \sectref{subsec:othergramm} investigates the interactions of apprehensionals in general with other grammatical features, especially grammatical person. \sectref{subsec:undesirability} draws attention to the varying degree of semanticisation for all apprehensional (sub)functions, while \sectref{subsec:modal} and \sectref{intro:functions} respectively explore the type(s) of modality instantiated by apprehensives, and their different pragmatic functions. \sectref{sec:diachrony} examines the diversity of the diachronic sources for apprehensionals. Finally, \sectref{sec:summaries} is a brief summary of the highlights of each contribution.


\subsection{Apprehensionality subtypes}
\label{subsec:apprsubtypes}

We use \textsc{apprehensional} as the cover term for grammatical markers or constructions of the entire functional domain, with the definition given in \sectref{subsec:appr}: grammatical markers or constructions which conventionally encode or pragmatically implicate that the situation described by the clause is an undesirable possibility. We refer to apprehensionals that conventionally encode undesirability as \textsc{dedicated}, and those that only pragmatically imply it as \textsc{non-dedicated}. Most authors in this volume adopt this definition and the terminological choices summarised in \tabref{tab1_alternative} and discussed in detail below.

\begin{table}
\caption{Terminology for the apprehensional constructions and their subtypes}
\label{tab1_alternative}
\begin{tabular}{L{1.4cm}L{4.6cm}L{4.6cm}}
\lsptoprule

& {\bfseries \textsc{dedicated} \textsc{apprehensional}}
 \textrm{(undesirability semantically encoded)} 
 & {\bfseries \textsc{non-dedicated} \textsc{apprehensional}}
 \textrm{(undesirability pragmatically implied)}\\\midrule

\textsc{main clause} & {\bfseries Apprehensive} & \textrm{\textbf{Possibility} \textbf{modal}} \\
& \hspace{8pt} \textrm{Specialised marker} & \\
& \hspace{8pt} \textrm{Transparent neg. optative} & \\
\textsc{adverb. clause} & \textrm{\textbf{Precautioning}} & \textrm{\textbf{Possible} \textbf{consequence}}\\
& \hspace{8pt} \textrm{\textit{Negative purpose}} & \\
& \hspace{16pt} \textrm{Specialised marker}    &   \\    
& \hspace{16pt} \textrm{Transparent negative} &  \\
& \hspace{16pt} \textrm{purpose} & \\
& \hspace{8pt} \textit{Possible undesirable} \\
& \hspace{8pt} \textit{consequence} \textrm{(`in-case')} & \\

\textsc{compl.\ clause} & \textrm{\textbf{`Fear'} \textbf{complementiser}} & \textrm{\textbf{General} \textbf{complementiser}} \\
\lspbottomrule
\end{tabular}
\end{table}

%\begin{table}
%\caption{Terminology for the apprehensional constructions and their subtypes}
%\label{tab1}
%\begin{tabularx}{\textwidth}{XX}
%\lsptoprule

%{\bfseries \textsc{Dedicated apprehensional}}

%\textrm{(undesirability semantically encoded)} 

%& {\bfseries \textsc{Non-dedicated apprehensional}}

%\textrm{(undesirability pragmatically implied)}\\\midrule

%\multicolumn{2}{c}{Main clause}\\\midrule
%{\bfseries Apprehensive}

%\hspace{8pt}  \textrm{Specialised marker} 

%\hspace{8pt}   \textrm{Transparent negative optative} & \textrm{\textbf{Possibility} \textbf{modal}}\\\midrule
%\multicolumn{2}{c}{Subordinated clause}\\\midrule
%\textrm{\textbf{Precautioning}} 

%\hspace{8pt}   \textrm{\textit{Negative purpose}}

%\hspace{16pt}   \textrm{Specialised marker}          

%\hspace{16pt}  \textrm{Transparent negative purpose}

%\hspace{8pt}   \textrm{\textit{Possible undesirable consequence}} 

%\hspace{8pt}  \textrm{(`in-case')} & \textrm{\textbf{Possible} \textbf{consequence}}\\

%\textrm{\textbf{`Fear'} \textbf{complementiser}} 
%& \textrm{\textbf{General} \textbf{complementiser}}\\
%\lspbottomrule
%\end{tabularx}
%\end{table}







A subdivision that has been influential for existing descriptions of apprehensional constructions is \citegen{lichtenberk1995}, who distinguishes between three broad construction types. The first type, which he terms \textsc{apprehensional}-\textsc{epi-stemic}, is typically realised as a construction at the level of the main clause. The other two are realised at the level of subordinate constructions. These in turn are subdivided into \textsc{precautioning adverbial clauses} and \textsc{complement clauses of `fear' predicates}. \citegen{lichtenberk1995} discussion centres on the \ili{Oceanic} language \ili{To'aba'ita}, where according to his description, the same marker occurs in all three functions, as illustrated in \REF{ex:intro4}.\footnote{We have slightly changed the glosses in some of the examples for the sake of clarity. Lichtenberk uses the gloss \textsc{lest} while we use \textsc{appr} here as a more general label for the apprehensional construction available in \ili{To'aba'ita}.}
\newpage

\ea 
\label{ex:intro4} 
\langinfo{To'aba'ita}{Austronesian: \ili{Oceanic}; Solomon Islands}{\citealt[296--8]{lichtenberk1995}} 
\ea \label{ex:intro4a} \textnormal{Apprehensional-epistemic} \\
\gll \textbf{Ada} wane `eri ka riki nau. \\
\textbf{\textsc{appr}} man that he:\textsc{seq} see me \\
\glt `[I fear] the man might see me.'  
\ex \label{ex:intro4b} \textnormal{Precautioning (negative purpose subfunction)} \\
\gll Nau ku	agwa `i	buira fau \textbf{ada} wane `eri ka riki nau.\\
I I:\textsc{fact} hide at behind rock \textbf{\textsc{appr}} man that he:\textsc{seq} see me\\ 
\glt `I hid behind the rock so that the man might not see me / lest the man see me.' 
  \newpage  \ex \label{ex:intro4c} \textnormal{Complement of `fear' predicate} \\
\gll Nau ku \textbf{ma'u} {`asia na'a} \textbf{ada} laalae to'a baa ki keka lae	mai keka thaungi kulu.\\
I I:\textsc{fact} \textbf{be.afraid} very \textbf{\textsc{appr}} later people that \textsc{pl} they:\textsc{seq} go hither they:\textsc{seq} kill us(\textsc{incl}) \\
\glt `I am scared the people might come and kill us.' / `I am scared lest the people should come and kill us.'  \\
\z
\z


 Languages like \ili{To'aba'ita}, which display a single apprehensional marker for the three semantically related functions outlined above, are not infrequent. However, languages exhibiting a clear morphological distinction between main clause and subordinate apprehensional constructions also exist. For example, \ili{Ese Ejja}, discussed in detail in \citet{vuillermet2018a}, has different markers for the (main clause) apprehensive and the (subordinate) precautioning function as illustrated in \REF{ex:intro5a}--\REF{ex:intro5b} below. In addition, \ili{Ese Ejja} also has a distinct apprehensional marker at the level of the noun phrase, i.e.\ a timitive (`for fear of') case marker/adposition, illustrated in \REF{ex:intro5c}.


\ea 
\label{ex:intro5}
\langinfo{Ese Ejja}{\ili{Pano-Tacanan}; Bolivia/Peru}{\citealt[270, 257]{vuillermet2018a}}\\
\ea \label{ex:intro5a} \textnormal{Apprehensive} \\
\gll A'a María wowi-jji, poki-\textbf{chana}! \\
\textsc{proh} María tell-\textsc{proh} go-\textsc{\textbf{appr}} \\
\glt `Don't tell Maria, she might go!'' \\
\ex \label{ex:intro5b} \textnormal{Precautioning} \\
\gll Owaya ekowijji shijja-ka-ani {\ob}\textbf{e}-jja-saja-ki \textbf{kwajejje}{\cb}. \\
\textsc{3erg} rifle clean-3\textsc{a-prs} \textsc{\textbf{prec}-mid}-block-\textsc{mid} \textsc{\textbf{prec}} \\
\glt `He is cleaning his rifle [lest it get blocked].' \\
    \ex \label{ex:intro5c} \textnormal{Timitive} \\
\gll Akwi iña-kwe iñawewa=\textbf{yajjajo}! \\
stick grab-\textsc{imp} dog=\textsc{\textbf{tim}} \\
\glt `Grab a stick for fear of the dog!'
\z
\z

 Timitives are known to show overlap with other apprehensional markers. For instance, \citet[74, 88]{dixon2009} mentions the frequent use of what he calls aversive cases as complementisers of `fear' predicates, and as precautioning subordinators in Australian languages. Timitives are excluded from the scope of this volume unless they show such formal overlap, as is the case e.g.\ in a number of \ili{Ngumpin-Yapa} languages (\citealt{chapters/browne}), in \ili{A'ingae} (aka.\ \ili{Cofán}; isolate; \citealt{chapters/dabkowski-anderbois}), and in \ili{Upper Tanana Dene} (\ili{Athabaskan-Eyak-Tlingit}; \citealt{chapters/lovick}). The formal overlap is motivated by the fact that a timitive-marked entity can often be interpreted as metonymically standing for an undesirable event (e.g.\ in \REF{ex:intro5c}: `for fear of the dog attacking').

In line with wide-spread usage (e.g.\ \cite{Dobrushina2006,pakendorfschalley2007,angeloschultzeberndt2016,vuillermet2018a}) we adopt the related term \textsc{apprehensive} for constructions that can occur at the main-clause level, i.e.\ \citegen{lichtenberk1995} apprehensional-epistemic use, corresponding to the \ili{Ese Ejja} construction illustrated in \REF{ex:intro5a}.\footnote{As will be discussed below (\sectref{intro:syntacticstatus}), the main clause status is not necessarily easy to establish, and there may be non-syntactic, though orthogonal, relevant criteria (like the necessary speaker's perspective) to distinguish apprehensive from precautioning uses.} Apprehensives mark the proposition in their scope as both possible and undesirable, though whether the possibilities considered are epistemic, as claimed by \citet{lichtenberk1995}, remains to be established rather than assumed, and may be language-specific. In some languages, apprehensives are part of a grammatical modal or mood paradigm (e.g.\  \ili{Ese Ejja}, see \REF{ex:intro5a}, or \ili{Kambaata} (Afro-Asiatic: \ili{Cushitic}; \citealt{chapters/treis})), while in others, they pertain to a different syntactic class. A negative marker may be part of a main-clause level construction to transparently indicate the speaker's desire for the situation described in the clause not to take place, in what can be described as a negative-optative apprehensional construction (see \citealt{chapters/wiemer} on \ili{Slavic} languages).

Where a language-specific construction pragmatically implicates rather than semantically encodes undesirability, a more general label, e.g.\ \textsc{possibility modal}, may be more appropriate.

For subordinate apprehensional clauses within a complex sentence construction, corresponding to the \ili{Ese Ejja} construction illustrated in \REF{ex:intro5b}, we will adopt  \citegen{lichtenberk1995} term \textsc{precautioning} clause. Like apprehensives, precautioning clauses convey both possibility and undesirability of the situation described, but require in addition the presence of a \textsc{preemptive} clause, which encodes an action to be taken or omitted in order to avoid the undesirable event expressed in the precautioning clause. (We deliberately chose \citeapo[264]{evans1995} term \textit{preemptive} over Lichtenberk's \textit{precautionary clause} because the latter is easily confused with \textit{precautioning}.) In these constructions, it is often a subordinator, e.g.\ \ili{English} \textit{lest}, that contributes the semantics of undesirable possibility. Preemptive clauses are frequently (according to \citet[24]{dixon2009}, ``most frequently'') directives, i.e.\ commands or prohibitions. Preemptive clauses are also frequently found to co-occur with apprehensives as defined above, which can contribute to blurring the boundary between main clause apprehensive and subordinate precautioning clauses (see \sectref{subsec:othergramm}).

Precautioning clause is in turn a cover term for \textsc{negative purpose} adverbial clauses (\citegen[298]{lichtenberk1995} \textit{avertive clauses}\footnote{We avoid the term \textit{avertive} for the negative purpose subtype (used e.g.\ by \cite{lichtenberk1995,schmidtkebode2009,vuillermet2018a}), since it has also been adopted for the unrelated function of expressing an action that is narrowly averted (e.g.\ \textit{almost}, \textit{be about to}) (e.g \cite{kuteva2000}, \cite{kutevaetal2019}). We retain Lichtenberk's term `\textit{in-case}' for the other clause type (an alternative is ``possible undesirable consequence'') despite the fact that the \ili{English} subordinator \textit{in case} is not restricted to undesirable events (see below).}) and `\textsc{in-case}' adverbial clauses, following \citegen[298]{lichtenberk1995}) term. Negative purpose clauses are characterised by having a direct causal link with the preemptive clause, and, consequently, the realisation of the preemptive action has an impact on the non-realisation of the undesirable event. For example, in \REF{ex:intro6}, not climbing up likely avoids the consequence of falling (from the wall or the tree) altogether. By contrast, a causal link is not required between an `in-case' precautioning clause and the preemptive clause: the undesirable event may happen anyway, but the preemptive action rather targets the (implicit) consequences of the event encoded by the `in-case' clause. In \REF{ex:intro7}, the undesirable event avoided by the realisation of the preemptive event is not the lightning strike itself but the potential harm to the climber. As shown by these examples, \textit{lest} is acceptable in both types of clauses, whereas the construction \textit{so that $\ldots$  not} transparently marks negative purpose only.

\ea \textnormal{Negative purpose precautioning clauses: causal link}
\label{ex:intro6} 
\ea\label{ex:intro6a} Don't climb up there \textbf{lest} you fall!
\ex\label{ex:intro6b} Don't climb up there \textbf{so that} you do\textbf{n't} fall!
\z
\z

\ea \textnormal{`In-case' precautioning clauses: no causal link}
\label{ex:intro7}
\ea[]{\label{ex:intro7a} She didn't climb up the tree \textbf{lest}/\textbf{in case} lightning struck.}
\ex[\#]{She didn't climb up the tree \textbf{so that} lightning would \textbf{not} strike.}\label{ex:intro7b} 
\z
\z

Note that the presence or absence of a causal relationship between the preemptive clause and the precautioning clause is related to the overall (lack of) preventability of the event: assuming that preventing lightning itself is impossible, taking precautions may at most avoid the undesirable potential consequences of this event. A number of contributions to this volume discuss the semantic and pragmatic relationship between negative purpose and `in-case' clauses for languages where both can be encoded by the same marker, e.g.\ \textit{pen} in \ili{Classical Hebrew} (Afro-Asiatic: \ili{Semitic}; \citealt{chapters/francez}), \textit{sa'ne} in \ili{A'ingae} (aka.\ \ili{Cofán}; isolate; \citealt{chapters/dabkowski-anderbois}), \textit{-lkut} in \ili{Archi} (\ili{East Caucasian}; \citealt{chapters/daniel}), \textit{tiple} in \ili{Mwotlap} (Austronesian: \ili{Oceanic}; \citealt{chapters/francois}); see also \sectref{subsec:prepreventability} in \citegen{chapters/vuillermet} questionnaire.


Negative purpose adverbial clauses, in turn, can be transparently marked as negative or employ a specialised marker like \ili{English} \textit{lest} \citep[130]{schmidtkebode2009}. If the subordinate clause type under investigation does not entail undesirability, a more general term than precautioning may be appropriate, e.g.\ \textsc{possible consequence} (cf.\ \cite[2]{dixon2009}) for clauses comparable to \textit{in case} clauses in \ili{English} \REF{ex:intro7a}.\footnote{Note that `in-case' clauses of possible consequence are semantically distinct from conditional clauses (pace \cite[136]{schmidtkebode2009}) -- \textit{in case} in \ili{English} covers both functions. The semantic distinction is formally reflected in other languages, cf.\ the two \ili{German} equivalents \textit{im Fall, dass} `if, in the case that' for conditionals and \textit{für den Fall, dass} `for the case that' for possible consequence clauses.}

Clauses serving as the complement of `fear' verbs can be described transparently as \textsc{`fear' complements} (the `fear' function in \citegen[297]{lichtenberk1995} terms). Here the undesirability reading of the subordinate clause is also conveyed by the semantics of the main clause predicate. The reason for including these in the apprehensional domain is that in many languages the same subordinator as in negative purpose or precautioning `in-case' clauses may be used, as illustrated for \ili{English} \textit{lest} in \REF{ex:intro8} and for \ili{To'aba'ita} \textit{ada} in \REF{ex:intro4c}. Very often a general complementation construction is used alongside or instead of an apprehensional construction in complements of `fear' verbs, as illustrated in \REF{ex:intro8} for \ili{English} and discussed for \ili{Ngumpin-Yapa} languages by \citet{chapters/browne}. 

\ea\label{ex:intro8} Tom was worried \textbf{lest} he fail / \textbf{that} he might fail the exam.
\z



\section{Morphosyntactic aspects}
\label{sec:morph}


\subsection{Syntactic status of the apprehensional clause}
\label{intro:syntacticstatus}

In some languages the syntactic status of the apprehensional clause is straightforwardly identified. Example \REF{ex:intro5} showed that three of the main apprehensional functions discussed in \sectref{subsec:apprsubtypes} align with clearly differentiated syntactic constructions and different morphological marking in \ili{Ese Ejja}: The Apprehensive\footnote{We capitalise the labels of individual language-specific morphemes.}   \nobreakdash-\textit{chana} is a main clause modal marker, the Precautioning \textit{e-$\ldots$kwajejje} is a subordinate construction, and the Timitive clitic \textit{=yajajjo} is a case marker \citep{vuillermet2018a}. However, for many languages the identification of the syntactic status of an apprehensional construction is a thorny issue, as already discussed in the literature (e.g.\ \cite[230]{Austin2021}). Language-specific descriptions often remain obscure about the reasons for identifying an apprehensional marker e.g.\ as a subordinate marker or mood marker. Moreover, a segmental marker of subordination may be absent (in which case the semantic relationship between the two clauses also remains underspecified). \citet[219]{verstraete2006b} emphasises that the absence of a relational marker does not exclude ``two clauses [to be] actually integrated via intonation, with intonational information taking over from segmental information as a relational marker.'' As a matter of fact, the role of intonation in creating clausal dependencies is typically underreported in grammatical descriptions.

One of the difficulties for distinguishing between main and subordinate apprehensional clauses more specifically is that their semantics frequently, possibly even typically, triggers the presence of another (preemptive) clause that encodes an action (to be) taken to avoid the undesirable consequence encoded or implied by the apprehensional clause. Thus, in language-specific discussions of apprehensionals, their subordinate status is often linked to the (near-)obligatory presence of a preemptive clause (e.g.\ \cite{smithdennis2021}). This is however not sufficient, as illustrated by the \ili{Kriol} example in \REF{ex:intro1}: While a preemptive clause is present, and the apprehensional marker \textit{bambai} has its origin in a temporal connective, \citet{angeloschultzeberndt2016} show that \ili{Kriol} apprehensional clauses have the status of main clauses, and that \textit{bambai} is a sentence connective in both its temporal and its apprehensive function. Thus \textit{bambai} serves to establish a pragmatic (rather than syntactic) dependency between the preemptive clause (if present) and the apprehensional clause, without being a subordinator (see \citealt{chapters/francois} for a similar argument for \ili{Mwotlap}). Similarly, the \ili{Ese Ejja} example \REF{ex:intro5a} shows that the mere presence of a preemptive clause cannot be a sufficient criterion for the subordinate status of the apprehensional clause. The \ili{Ese Ejja} Apprehensive \nobreakdash-\textit{chana} is syntactically a main clause marker for two reasons: it belongs to the obligatory main clause tense/mood paradigm, and only the main clause pronoun sets are available with this paradigm \citep[267--268]{vuillermet2018a}.

The appearance of subordination may also arise from the use of a main clause apprehensional construction (i.e.\ an apprehensive) in a reported speech\slash thought construction, as in `She did not go to the party, saying ``I might get bored'''. This construction is illustrated in \REF{ex:intro9} for \ili{Tariana}. Similar examples of reported speech constructions are discussed by  \citet{chapters/daniel} for \ili{Mehweb} (aka.\ North Central \ili{Dargwa}, \ili{East Caucasian)},  \citet{chapters/rhee-kuteva} for a number of the \ili{Korean} apprehensional constructions discussed in their paper, and \citet[1170]{jacques2021} for \ili{Japhug} (\ili{Sino-Tibetan}: \ili{Gyalrongic}).

\ea \label{ex:intro9} 
\langinfo{Tariana}{\ili{Arawakan}; Brazil/Colombia}{\citealt[384]{aikhenvald2003}}\\
\gll {\ob}Nesu-\textbf{da} na-swe di-\textbf{a}-\textbf{ka}{\cb} di-piti-naka. \\
\textsc{3pl}$+$shit-\textsc{\textbf{vis.appr}} \textsc{3pl-}stay$+$\textsc{caus} \textsc{3sg.nf-}\textbf{say-\textsc{sub}} \textsc{3sg.nf-}drive.away-\textsc{prs.vis} \\
\glt `He is driving (the hens) away for fear they might shit.' (lit.\  `$\ldots$ saying they might shit.')
\z


 Even apparent complements of `fear' verbs may in reality be syntactically independent clauses. According to \citet[619]{tsunoda2011}, a \ili{Warrongo} speaker translated the \ili{English} prompt `The woman is afraid that her child might be drowned' with the utterance in \REF{ex:intro10} (glosses adjusted). 

\ea \label{ex:intro10} 
\langinfo{Warrongo}{\ili{Pama-Nyungan}; Australia}{\citealt[619]{tsunoda2011}}\\
\gll Warrngo wanbali-n	galbin jojorri-ngga. \\
woman fear-\textsc{nfut} child be.drowned-\textsc{appr} \\
\glt `The woman is afraid -- the child might be drowned.'
\z

\noindent The translation that we provided with \REF{ex:intro10} reflects the observations that the verb \textit{wanbali-} `fear, be afraid' can occur in an intransitive clause without a complement \citep[401]{tsunoda2011}, that the \ili{Warrongo} Apprehensive marker is found in independent main clauses \citep[286--287]{tsunoda2011}, and that there is no marker of subordination in such complex constructions. 


In other cases, there may be genuine arguments for claiming that an apprehensional construction has subordinate status while not showing many signs of deranking (i.e.\ formal correlates of subordination),\footnote{\citet[439]{schmidtkebode2012} observes that cross-linguistically, subordinate precautioning/negative purpose clauses show fewer hallmarks of de-ranking than positive purpose clauses.} or that the same construction can alternatively have main clause and subordinate clause status, as discussed in some detail by \citet{lichtenberk1995} and \citet{Dobrushina2006}. This syntactic ambiguity can come about through syntactic tightening of two juxtaposed clauses (see \citealt{chapters/wiemer}). Conversely, a main clause use of an apprehensive marker could arise through insubordination of a subordinate precautioning clause (an incipient case is arguably the \ili{English} expression \textit{lest we forget}), with potentially formal remnants of such an origin (the subjunctive in the case of \textit{lest}). For example, the Apprehensive in \ili{To'aba'ita} illustrated in \REF{ex:intro4a} does not belong to the same paradigm as other moods and still displays a subject/tense marker from the sequential paradigm, typical of subordinate clauses \citep[296]{lichtenberk1995}. 

Apart from general language-specific correlates of subordination, one potential semantic criterion for disentangling the two syntactic levels is the judge of undesirability: The judge in main clause apprehensionals is always the speaker, the speaker-addressee dyad, or another contextually determined judge (\cite[164--170]{green1989},  \citealt{chapters/fedotov}; \citealt{chapters/francois}, \cite[272--273]{vuillermet2018a}, \citealt{chapters/vuillermet}: \S2.3).\footnote {See also the speaker's vs.\ agent's perspective in \citegen[68--70]{verstraete2006a} discussion of mood use in past intentional constructions. \citet[68--69]{verstraete2006a} also discusses grammaticalised ``agent-binding'' morphemes, which bind the modal feature (e.g.\ apprehensive) to the agent of the main clause (e.g.\ yielding a precautioning clause), in the Australian languages \ili{Nyikina} (\ili{Nyulnyulan}) and \ili{Ngiyambaa} (\ili{Pama-Nyungan}).} By contrast, in a precautioning clause, the judge has to be the subject of the main clause. In \REF{ex:intro3}, the judge is the father who tied his daughter's laces, and in \REF{ex:intro5b}, it is the hunter cleaning his rifle. One exception to the generalisation that the judge of undesirability in a precautioning clause has to be the subject of the main clause arises when this main clause is a directive. In \textit{\textbf{Don't swim}  in this lake, lest the crocodiles eat you!} or \textit{Put your jacket on, lest you catch a cold!}, the judge is typically the speaker who utters the directive.

As already mentioned above, if a language does not have a precautioning construction, the functional need for assigning the judge role to someone other than the speaker can be met in languages like \ili{Tariana} by using an apprehensive clause in a reported speech construction as in \REF{ex:intro9}: this binds the judgement of undesirability to the main clause agent whose thoughts or speech are reported. 

In some of this volume's contributions, the syntactic status of the apprehensional clause can straightforwardly be allocated to a single construction type. Many of the contributions, however, are dealing with the use of the same construction in what appear to be both apprehensive and precautioning clauses. Some of these languages transparently use the reported speech strategy mentioned above, while others seem to have developed the independent from the dependent construction via an insubordination pathway (\citeauthor{chapters/dabkowski-anderbois} for \ili{A'ingae}; \citeauthor{chapters/daniel}  for \ili{Archi}, \citeauthor{chapters/francez} for \ili{Hebrew}). On the other hand, \citeauthor{chapters/wiemer} cautions against the assumption of insubordination for \ili{Slavic} languages, and \citeauthor{chapters/francois}  (cf.\ also \cite{francois2003}) considers apprehensional clauses in \ili{Mwotlap} to have a pragmatically rather than syntactically dependent status.




\subsection{Apprehensionals and their interaction with other grammatical domains}
\label{subsec:othergramm}

Another fertile area of investigation is the interaction of apprehensional constructions with other grammatical domains. We will briefly discuss the interaction with grammatical person, tense/aspect, and negation. Interaction with grammatical person concerns the restriction of apprehensional markers to specific grammatical persons, or to specific combinations of grammatical persons \citep{vuillermet2018b}. (Grammatical persons may also more or less directly impact the semantic interpretation of  apprehensives, which is dealt with in \sectref{subsec:modal}.) For instance, the \ili{Andoke} (isolate; Colombia) Apprehensive is only available with 2\textsuperscript{nd} person \citep[89]{landaburu1979}, while the \ili{Harakmbut} (aka \ili{Amarakaeri}, isolate; Peru) Apprehensive seems to only be available with 1\textsuperscript{st} and 3\textsuperscript{rd} person (\cite[240]{helbergchavez1989}, \cite[222]{tripp1995}, An Van linden, p.c.\  August 2019). \ili{Matsés} displays another type of restriction, in that it has several Apprehensive markers specialised for distinct persons. Among them are \textit{-mane} \REF{ex:intro11a}, which exclusively occurs with 1\textsuperscript{st} person subjects, and \textit{-nunda} \REF{ex:intro11b}, with third person subjects.


\ea \label{ex:intro11} 
\langinfo{Matsés}{\ili{Pano-Tacanan}; Brazil/Peru}{\citealt[439--440]{fleck2003}}\\
\ea \label{ex:intro11a} \textnormal{1\textsuperscript{st} person subject Apprehensive marker} \\
\gll Mibi 	cues-mane. \\
2:\textsc{abs} hit-\textsc{\textbf{appr:1}} \\
\glt `I might accidentally hit you.' \\
$[$e.g.\  if you stand behind me as I split firewood] \\
\ex \label{ex:intro11b} \textnormal{3\textsuperscript{rd} person subject Apprehensive marker} \\
\textit{chononda} \\
\gll cho-\textbf{nunda} \\
come-\textsc{\textbf{appr:3($>$3)}} \\
\glt `(he) might come' \\
$[$e.g.\ ``don't go to your field, your lover might come'', or ``don't go there with your 1\textsuperscript{st} lover, the other lover might come.''] (D.\ Fleck p.c., September 2017)]
\z
\z

 Another two \ili{Matsés} Apprehensive markers contrast in having different grammatical persons as objects: \textit{-panondac}, illustrated in \REF{ex:intro12}, only occurs with 1\textsuperscript{st} and 2\textsuperscript{nd} person objects, and \textit{-nunda} with 3\textsuperscript{rd} person objects (see also \REF{ex:intro11b}).

\newpage
\ea \label{ex:intro12} 
\langinfo{Matsés}{\ili{Pano-Tacanan}; Brazil/Peru}{\citealt[440]{fleck2003}}\\
\ea \label{ex:intro12a} \textnormal{3$>$1 or 2 Apprehensive marker} \\
\gll Nisi-n pe-\textbf{panondac}. \\
snake-\textsc{erg} bite-\textsc{\textbf{appr:3$>$1-2}} \\
\glt `Now maybe a snake will bite me.' \\
$[$Said, e.g.\  as the speaker is going off to hunt on a rainy day.] \\
\ex \label{ex:intro12b} \textnormal{3$>$3 Apprehensive marker} \\
\gll Nisi-n pe-\textbf{nunda}. \\
snake-\textsc{erg} bite-\textsc{appr:3($>$3)} \\
\glt `(Be careful), a snake might bite (e.g.\  our son).'
\\
Not: `$\ldots$ bite you/me.'
\z
\z

\ili{Ye'kwana} (\ili{Cariban}; Venezuela; \cite[232--233]{caceres}) is another language whose apprehensive construction has two forms depending on the grammatical person combination: the negative marker \textit{an-} is necessarily present if a third person patient is involved, while it is absent in all other configurations.

Apprehensionals may also differ with respect to how they interact with tense or aspect. The impact of aspect on the availability of the prohibitive vs.\ preventive reading is well documented in \ili{Slavic} languages (see \sectref{sec:diachrony} and (\citealt{chapters/wiemer})). In \ili{Ese Ejja}, no tense marker can occur together with the Apprehensive since they belong to the same paradigmatic slot (\cite[367, 449]{vuillermet2012}; \cite[268]{vuillermet2018a}). In \ili{Hup} (\ili{Naduhup}; \cite[607--608]{epps2008}), the Future marker \nobreakdash-\textit{teg \~{ } \nobreakdash-te-}, which indicates a certain future event, can also not co-occur with the Apprehensive (nor with the Imperative, the Negative, the Habitual and some conditional expressions). However, the Future Contrast marker \textit{tán}, which indicates a ``potential or hypothetical future'', can (see \REF{ex:intro22} further below). Here it seems that possibility of co-occurrence is determined by the semantics of the two Future markers and that of the Apprehensive. 

As already stated in \sectref{introintro}, some apprehensional constructions (negative optatives, negative purpose clauses) include a transparent marker of negation (\citealt{chapters/wiemer}). Other apprehensional constructions, on the other hand, may include expletive negation, that is, a negation marker which does not transparently negate the clause it occurs in, but rather emphasises the undesirability of the potential situation. Expletive negation is quite common in the complement of `fear' predicates (\cite{jinkoenig2020, dobrushina2021}). Olguín Martinéz (\citeyear[601--602]{olguinmartinez2023a}; \citeyear[16]{olguinmartinez2023b}) also finds that the presence of negation in precedence (`before') clauses -- in languages where it is optionally possible -- tends to mark these clauses as non-factive and apprehensional, as opposed to factive clauses with a purely temporal reading. Expletive negation in \ili{French} in a `fear' complement and a precedence construction is illustrated in \REF{ex:intro13a} and \REF{ex:intro13b}, respectively. In both constructions, the (optional) negative marker \textit{ne} does not negate the undesirable event; the full \ili{French} negative construction \textit{(ne) $\ldots$ pas} would be necessary to achieve this.

\ea 
\label{ex:intro13a}
\langinfo{French}{Indo-European: \ili{Italic}}{\citealt[39]{jinkoenig2020}} \\
\gll J'ai peur qu'il {\op}\textbf{ne}{\cp} pleuve demain. \\
I.have fear that.it \textsc{\textbf{expl.neg}} rain:\textsc{3sg.sbjv} tomorrow\\
\glt `I fear that it will rain tomorrow.'\\
Not: `I fear that it will \textbf{not} rain tomorrow.'\\
\z

\ea 
\label{ex:intro13b}
\langinfo{French}{Indo-European: \ili{Italic}}{\citealt{tahar2021}} \\
\gll Jim doit attraper le vase grec avant qu'il  {\op}\textbf{ne}{\cp} tombe.\\
Jim must catch the vase greek before that.it \textsc{\textbf{expl.neg}} fall:-\textsc{3sg.sbjv} \\
\glt `Jim must catch the Greek vase before it falls.'\\
Not: `$\ldots$ before it does \textbf{not} fall.'
\z


 Whether a formally negative apprehensional construction transparently negates the clause it occurs in or is to be interpreted as expletive negation may depend on the linguistic context. For example, \citet{chapters/fedotov} argues that the negation marker in \ili{Gban} constitutes transparent negation in negative purpose constructions, but is expletive in all other apprehensional uses (see also \citetv{chapters/francez}).

Finally, there are apprehensionals which cannot occur with negation at all (e.g.\ in \ili{Archi} and \ili{Akhvakh}, \citetv{chapters/daniel} -- or more precisely, they cannot change polarity if they already include an expletive or transparent negation marker. In \ili{Upper Tanana Dene}, precautionary clauses are only compatible with affirmative main clauses (\citealt{chapters/lovick}).



\section{Semantics of the apprehensional marker}
\label{sec:semantics}

\subsection{Degree of semanticisation of the undesirability component}
\label{subsec:undesirability}


A major issue to consider in the domain of apprehensionality is whether the construction in question semantically encodes undesirability (\textit{dedicated apprehensional}, cf.\ \sectref{subsec:apprsubtypes}), or whether this is merely pragmatically implied in certain contexts (\textit{non-dedicated apprehensional}), for example as a result of its frequent use in warnings and other typical apprehensional contexts. Diachronically, semanticisation of a pragmatic implicature may of course lead from non-dedicated to dedicated apprehensionals.

The distinction between dedicated apprehensionals and more general constructions which are also frequently used in apprehensional contexts can be illustrated with the contrast between \textit{in case} and \textit{lest} in \ili{English}. \textit{Lest} is a subordinator specialised for undesirability \citep{lópezcouso2007}, while \textit{in case} has a general meaning of possible consequence (\cite[23--25]{dixon2009}; see also \sectref{subsec:apprsubtypes}). Although often used in contexts of undesirability, \textit{in case}, but not \textit{lest}, can introduce subordinate clauses that describe desirable events, as illustrated in \REF{ex:intro14}.

\ea\label{ex:intro14} You should keep your lottery ticket in a safe place, \textbf{in case}/\#\textbf{lest} you win.
\z

 An example from the modal domain for an apprehensional strategy that relies on pragmatic inference comes from the Australian language \ili{Jaminjung}-\ili{Ngaliwurru} \citep{schultzeberndtcaudal2019}. In affirmative clauses, the irrealis modal prefix (which contrasts with a second modal prefix used in desiderative and deontic contexts) is overwhelmingly used in warnings such as \REF{ex:intro15a}, and in most contexts implies an undesirable possibility even in the absence of a preemptive clause. However, examples such as \REF{ex:intro15b} show that it is still a modal of general possibility which does not semantically encode undesirability\footnote{Additional evidence comes from the fact that the \textsc{irr} modal is the only modal appearing in negative contexts, i.e.\ it expresses general impossibility without an undesirability component.} (cf.\ also \citet{chapters/browne} for a similar pattern in some neighbouring Australian languages).

\ea 
\label{ex:intro15} 
\langinfo{Jaminjung-Ngaliwurru}{\ili{Mirndi}; Australia}{fieldwork Schultze-Berndt}\\
\ea \label{ex:intro15a}
\gll Gurrany mard ya-nth-angu, guyug-burru, \textbf{ya}-nth-irna! \\
\textsc{neg} touch \textsc{irr-2sg$>$3sg-}get/handle fire-\textsc{propr} \textsc{\textbf{irr-}2sg-}burn \\
\glt `Don't touch it, it is hot like fire, you might get burned!' \\
\ex \label{ex:intro15b}
\gll Ning \textbf{ya}-nth-ina lidburrg-ni. \\
break.off \textsc{\textbf{irr}-2sg$>$3sg-}chop axe-\textsc{inst} \\
\glt (Echidnas are found under rocks.) `You might/can kill one with an axe.'
\z
\z

If a marker still conveys undesirability in pragmatic contexts that are typically considered desirable (or else is judged unacceptable in such contexts), it can be concluded that it conventionally codes this meaning component and is thus a dedicated apprehensional marker. An example is the \ili{Ese Ejja} Apprehensive \nobreakdash-\textit{chana}. Even if it occurs with a typically positive event like the verb \textit{swa-} `laugh', as in \REF{ex:intro16}, the event of laughing can only be interpreted in the sense of a mockery.


\ea 
\label{ex:intro16}
\langinfo{Ese Ejja}{\ili{Pano-Tacanan}; Bolivia/Peru}{fieldwork Vuillermet} \\
\gll Marina swa-\textbf{chana}, mi='ba-majje! \\
Marina laugh-\textsc{\textbf{appr}} \textsc{2sg.abs=}see-\textsc{tmp.ss} \\
\glt `Marina might laugh at seeing you!'
\z

 Another example is the \ili{German} connective \textit{nachher}, which in its original temporal sense means `afterwards, later', but has in colloquial \ili{German} acquired a second, apprehensional sense. In this sense, it is a dedicated apprehensional marker according to this test. Thus, the apprehensional clause in \REF{ex:intro17} can only be interpreted as describing an undesirable situation; combining it with an expression of positive evaluation as in \REF{ex:intro17b} leads to infelicity.

\ea \label{ex:intro17} 
\langinfo{German}{Indo-European: \ili{Germanic}}{\cite{fallerschultzeberndt2018}}\\
\ea \label{ex:intro17a}
\gll Ich glaube, ich nehme lieber nicht am Gewinn-spiel	teil. \textbf{Nachher} gewinne ich noch \ldots. \\
I think:\textsc{1sg.prs} I take:\textsc{1sg.prs} better not at:\textsc{def} win-game part \textsc{\textbf{appr/}}\textbf{later} win:\textsc{1sg.prs} I \textsc{add} \\
\glt `I think I better not take part in this prize draw. I might actually win [a trip] \ldots.' \\
Context: $\ldots$ and would not be able to watch the Champions League final at home.\footnote{\url{http://www.hifi-forum.de/viewthread-144-7707-4.html} (last accessed 2024-03-18)} \\
\ex[\#]{\label{ex:intro17b}
\gll \textbf{Oh} \textbf{super!} \textbf{Nachher} gewinne ich noch \dots \\
oh great \textsc{\textbf{appr/}}\textbf{later} win:\textsc{1sg.prs} I \textsc{add} \\
\glt Intended: `Great! I might actually win \dots.'}
\z
\z

 If undesirability has been established as a semantic, that is, conventional meaning component of the construction, a further question arises regarding the type of semantic meaning. Is it part of the at-issue meaning, that is, the main point of the utterance, or is it backgrounded in some way? One test to distinguish at-issue from not-at-issue meaning  relies on the observation that only at-issue meaning is accessible to direct expressions of assent/dissent (\cite{faller2002,papafragou2006,tonhauser2012,murray2021}). As shown in \REF{ex:intro18}, this test establishes that the proposition denoted by the apprehensive clause in \REF{ex:intro18} is at-issue, while its undesirability is not. In contrast, with \textit{I fear} in \REF{ex:intro19}, both the proposition denoted by the complement clause and its undesirability may be at-issue.
\newpage
\ea \label{ex:intro18} 
\langinfo{German}{Indo-European: \ili{Germanic}}{\cite{fallerschultzeberndt2018}}\\\exi{A\phantom{'}:}
\gll Spiel nicht mit der Vase, \textbf{nachher} geht sie noch kaputt. \\
play:\textsc{imp} \textsc{neg} with \textsc{def.dat.fem} vase \textsc{\textbf{appr/}}\textbf{later} go:\textsc{3sg.prs} \textsc{3sg.fem} \textsc{add} broken \\
\glt `Don't play with the vase, it might break (and that would be undesirable).'
\exi{B\phantom{'}:} `No, it won't break.' ($\Rightarrow$ breaking at-issue)
\exi{B':}[\#]{`No, that wouldn't be a bad thing -- it's ugly anyway.' ($\Rightarrow$ undesirability not at-issue)}
\z


\ea \label{ex:intro19} \textnormal{English}
\exi{A\phantom{'}:} \textbf{I fear} the vase might break.
\exi{B\phantom{'}:} No, it won't break. (\textnormal{$\Rightarrow$ breaking at-issue})
\exi{B':} (Don't worry,) that wouldn't be a bad thing -- it's ugly anyway. (\textnormal{$\Rightarrow$ undesirability at-issue})
\z


\noindent \citet{chapters/francez} also claims for the \ili{Hebrew} apprehensionals that the undesirability component is not at-issue and specifically analyses it as a conventional implicature.


\subsection{Apprehensives as a type of modal}
\label{subsec:modal}

While all apprehensionals have modal semantics as markers of possibility and undesirability, apprehensives, that is those markers that can occur in main clauses, in many languages are part of the modal system, and specifically of the temporal/modal/mood paradigm (see, in this volume, \citeauthor{chapters/browne} for a subset of \ili{Ngumpin-Yapa} languages, \citeauthor{chapters/fedotov} for \ili{Gban},  \citeauthor{chapters/francez} for one of the \ili{Hebrew} apprehensionals, \citeauthor{chapters/francois} for \ili{Mwotlap}, \citeauthor{chapters/treis} for \ili{Kambaata}, and also \citealt{vuillermet2012} \citeyear{vuillermet2018a} for \ili{Ese Ejja}). That apprehensives are a subtype of modal is also supported by a diachronic relationship between possibility modals of various types and apprehensives (see further \sectref{sec:diachrony}).

A careful analysis of such markers regarding the paradigmatic oppositions and/or syntagmatic co-occurrence with other modals or mood will go towards elucidating their semantic contribution to the clause. So far, it is an unresolved question what subtype(s) of modality can be instantiated by apprehensive modals. They are described as a mixed modal category with both epistemic and attitudinal components (potentiality and undesirability) by \citet[293]{lichtenberk1995}, but it remains to be established whether apprehensives are always truly epistemic. For example, apprehensives may be used to mark general undesirable dispositions such as \textit{Dogs may bite} or statistical undesirable likelihoods such as \textit{Your lottery ticket may/will likely not win}, and these are arguably not epistemic. It is also conceivable that an apprehensive could be used to mark obligations as undesirable. A diagnostic for potential epistemic status is interaction with tense. It has been observed that non-epistemic modals are necessarily future-oriented (\cite{condoravdi2002,rullmannmatthewson2018}), that is, the state of affairs described must hold after the relevant circumstances are established (the circumstances themselves can be established in the past, present or future). Epistemic modals, however, can also be past- or present-oriented, that is, the possible state of affairs may hold not only after but also before or at the time at which the epistemic agent's relevant beliefs hold. For example, the speaker of \REF{ex:intro20} considers it a possibility that their oral exam date is in the past based on their current knowledge (other students have already had their oral exam). This suggests that \ili{German} \textit{nachher} is indeed used epistemically (similarly for \ili{Hebrew}, \citealt{chapters/francez}).

\ea \label{ex:intro20} \textnormal{German} \\
\gll {Ich will auch meinen mündlichen termin!!!!  nachher war er schon und ich bin nicht informiert worden!}\\
\textsc{1sg} want.\textsc{1sg.prs} also my oral date 
\textsc{\textbf{appr/}}\textbf{later} be.\textsc{3sg.pst} \textsc{3sg} already and \textsc{1sg} \textsc{aux.1sg} \textsc{neg} informed \textsc{pass.pst} \\
\glt `I want the date for my oral [exam] too! It may have taken place already and I wasn't informed!'\footnote{\url{https://www.medi-learn.de/foren/archive/index.php/t-43472-p-66.html} (last accessed 2024-03-15).}
\z

 However, some Apprehensives, e.g.\  in \ili{Gban} \citep{chapters/fedotov}, in \ili{Akhvakh} \citep{chapters/daniel}, and in \ili{Kambaata} \citep{chapters/treis}, cannot occur with the past at all, which indicates that they cannot be past-oriented, and, consequently, that they may be non-epistemic.

\subsection{Functions}
\label{intro:functions}

Since apprehensives are strongly associated with contexts of warnings and associated directives, such utterances are often performative. \citet[179]{bybeeetal1994} therefore subsume apprehensives (their ``admonitives'') under ``speaker-oriented modality'' (together with imperatives, optatives, etc.). Cross-linguistic evidence for the performative nature of apprehensive modals comes from their attested diachronic relationship with prohibitives (see below for discussion and references) and optatives (\citealt{Dobrushina2006}, \citealt{chapters/lahaussois}, \citealt{chapters/wiemer}), and their occurrence in the same paradigmatic slot as the imperative, e.g.\ in \ili{Kotiria}/\ili{Wanano} (\ili{Tucanoan}; Brazil and Colombia; \cite[268, 309]{stenzel2013}). \citet{chapters/treis} discusses to what extent the Apprehensive paradigm in \ili{Kambaata} may be considered a subtype of the directive mood, as it shares some semantic features with negative imperatives and jussives, but also considers the possibilities that it is a subtype of the indicative mood or constitutes its own category. 

In this context, an interesting but underexplored semantic parameter for apprehensives is whether they can occur in both warnings and threats, or only warnings; in other words, whether or not they are restricted to unintentional undesirable events (\cite{vuillermet2017b}, \cite[273--274]{vuillermet2018a}, \cite{VuillermetM2018questionnaire}). For example, \ili{Matsés} is a language where only the warning reading, illustrated in \REF{ex:intro21}, is available, not the threat reading (D.\ Fleck p.c., September 2017).\footnote{\citet{fleck2003} calls the \ili{Matsés} Apprehensive a \textit{Future Potential}. The number following the adapted gloss \textsc{appr} stands for the person of the subject, namely `1\textsuperscript{st} person' in example \REF{ex:intro21}.}

\ea \label{ex:intro21} 
\langinfo{Matsés}{\ili{Pano-Tacanan}; Brazil/Peru}{\citealt[439]{fleck2003}}\\
\gll Mibi cues-\textbf{mane}. \\
2:\textsc{abs} hit-\textsc{\textbf{appr:1}} \\
\glt `I might accidentally hit you.' [e.g.\  if you stand behind me as I split firewood] -- not: `Watch out I'm going to hit you!'
\z


In languages where it is available, the threat reading seems to be favoured by a number of factors. One is grammatical person: \citet[631]{epps2008} observes that a ``threat is [\ldots] the default interpretation when the subject is in the first person'' in her \ili{Hup} corpus, as illustrated in \REF{ex:intro22}.

\ea \label{ex:intro22}
\langinfo{Hup}{\ili{Naduhup}; Brazil/Colombia}{\citealt[231]{epps2008}}\\
\gll ʔám-ǎn ʔãh yɔmɔ̌y yók tán-ãh! \\
\textsc{2sg-obj} \textsc{1sg} anus stab.\textsc{appr} \textsc{fut.cntr-decl} \\
\glt `I'll stab you in the anus!' 
\z


\citet[273--4]{vuillermet2018a} suggests that the feature of control in the verb semantics is another factor: in \textit{(Watch out) I might fall on you!}, the typical reading without context is a warning rather than a threat. These two factors (1\textsuperscript{st} person subject and $+$control verb) are typical correlates of the necessary broader conditions that must hold for a threat reading (\cite{searle-taxonomy,salgueiro2010}), but these can also be met in other ways. Thus, \textit{He might burn your house!} is a threat if the speaker commits to bringing about the event, which requires them to have direct or indirect control over it, and the event is (or believed to be) undesirable primarily to the addressee (rather than to the speaker or to both). Often, threats are conditional on the addressee fulfilling a request, for example, handing over money, that is, the addressee is given a way of avoiding the undesirable event or its consequences (see also \citealt{chapters/fedotov}).

A prohibitive reading is also frequently available with second person apprehensives, typically with events that the addressee can control, as has been discussed in the literature (\cite[57]{Dobrushina2006}, \cite{pakendorfschalley2007}, see also \citealt{chapters/wiemer}). In \ili{Wintu}, the 2\textsuperscript{nd} person Apprehensive, which originally comes from an optative `may, might' \citep[193]{pitkin1984}, receives an apprehensive or a prohibitive interpretation, depending on the degree of control the addressee has over the event: getting drowned or getting hurt typically occur unwillingly (lack of control), resulting in a warning interpretation of \REF{ex:intro23a}, while telling a lie is typically controllable, and therefore more likely intended as a prohibitive as in \REF{ex:intro23b}.

\ea \label{ex:intro23} 
\langinfo{Wintu}{\ili{Wintuan}; United States}{\citealt[115]{pitkin1984}}\\
\ea \label{ex:intro23a} \textnormal{Apprehensive reading} \\
\textit{tuku\textbf{ken}}  \textnormal{`You might get drowned, watch out!'}\\
\textit{podo'lu\textbf{ken}}  \textnormal{`Look out, you might get hurt!'}
\ex \label{ex:intro23b} \textnormal{Prohibitive reading}\\
\textit{bala\textbf{ken}} \textnormal{`Don't tell a lie!'}
\z
\z


\noindent More generally, the preventability of the undesirable event plays a role in determining the function of apprehensives (\citealt{chapters/fedotov}, \citealt{chapters/vuillermet}). This semantic variation is discussed further in the next section on diachrony.



\section{Diachrony}
\label{sec:diachrony}

\citet{lichtenberk1995} makes the strong claim that apprehensive main clause markers have their origin in precautioning clauses via the function of complement of `fear' predicates (assuming a pathway of insubordination):

\begin{quote}
    The implication of a situation being undesirable, or indeed feared, becomes explicit by virtue of clauses encoding such situations as subordinate to predicates of fearing. [\ldots] Over time, through metonymy based on cooccurrence of linguistic forms, the notion of apprehension comes to be associated with the \textsc{lest} element, and the predicates of fearing become expendable. \citep[319]{lichtenberk1995}
\end{quote}

 However, \citet[303]{lichtenberk1995} also notes that the Toba'a'ita apprehensional marker \textit{ada} has a likely origin in a verb of perception used as an attention getter (`look (out)!'). Attention getters are indeed attested in other languages as source constructions for apprehensionals, e.g.\ (possibly) in \ili{Akhvakh} (\citealt{chapters/daniel}); see also the discussion of \ili{Spanish} \textit{ten cuidado} `be careful' as the source for the apprehensional \textit{kuyra} in \ili{Huallaga Huánuco Quechua} (see \REF{ex:intro26}). An attention getter as grammaticalisation source however suggests a main clause origin of the apprehensive construction rather than insubordination via a complement of `fear'. 

Subsequent studies have shown a wide range of possible source constructions for apprehensional constructions; a number of them are discussed by the contributors to this volume, to the extent that the data for a given language permit conclusions about the origin of the construction. For example, \citet{chapters/rhee-kuteva} report, for \ili{Korean}, the grammaticalisation of `fear' predicates (rather than insubordination of a `fear' complement clause) as one possible pathway for the semanticisation of undesirability.

The link between apprehensionals and prohibitives -- already alluded to and illustrated with examples from \ili{Wintu} in \sectref{intro:functions} -- is well attested; it has been discussed in a typological perspective by \citet[50--63]{Dobrushina2006} and \citet{pakendorfschalley2007}, and is confirmed by a number of contributions to this volume (\citeauthor{chapters/lahaussois}, \citeauthor{chapters/treis}, \citeauthor{chapters/wiemer}) and elsewhere \citep{smithdennis2021}. \citet{pakendorfschalley2007} posit a pathway from general possibility marker to prohibitive via an apprehensive modal: A sentence describing a potential undesirable event can be interpreted as a warning (as in \REF{ex:intro25a}) and therefore be reinterpreted as a directive to avoid the undesirable event \REF{ex:intro25b}.


\ea\label{ex:intro25}
\ea\label{ex:intro25a} {\op}Be careful,{\cp} you might slip!
\ex\label{ex:intro25b} {\op}Be careful,{\cp} don't slip!
\z
\z

 \citet[62]{Dobrushina2006} considers both this and the reverse pathway plausible. As pointed out by her (\citeyear[61--63]{Dobrushina2006}), and discussed in some detail by \citet{chapters/wiemer}, an important link between the apprehensive and prohibitive function are a subtype of prohibitives termed \textit{preventives}, prominent in the literature on \ili{Slavic} languages since here they are distinguished from canonical prohibitives in terms of aspectual marking (perfective vs.\ imperfective) (see also e.g.\ \cite{Xrakovskij1990,kuehnast2008}). The hallmark of preventives is that they are formally prohibitives but encode an (undesirable) event that is not directly controllable by the addressee, as in \REF{ex:intro25b}. Rather, their pragmatic function is to invite the addressee to take measures to prevent the event in question -- in other words, they warn of a potential negative consequence rather than directly prohibiting the action leading to it. In either case, pragmatic enrichment is at work: a prohibitive can be pragmatically employed to encode a warning (if avoiding an event is in the interest of the addressee rather than another party), as in \REF{ex:intro25b}, leading to a re-analysis of the prohibitive as an apprehensive. The reverse could ultimately lead to a re-analysis of the apprehensive as a prohibitive. 


The pathway from a more general possibility modal (which itself will have ultimate lexical sources such as motion verbs or expressions of uncertainty) to an apprehensive modal, via pragmatic strengthening of the undesirability component (\cite[211]{bybeeetal1994}; \cite[34--36]{Dobrushina2006}; \cite{pakendorfschalley2007}), was already briefly discussed above and in \sectref{subsec:undesirability} and \sectref{subsec:modal}. Several contributions to this volume examine the degree of semanticisation of the undesirability component for apprehensive strategies originating in lexical verbs of cogitation via markers of epistemic possibility\slash uncertainty (\citeauthor{chapters/rhee-kuteva} for \ili{Korean}), in a modal adjective encoding propensity (\citeauthor{chapters/francez} for Modern \ili{Hebrew}), in a general possibility marker (\citeauthor{chapters/browne} for various Ngumpin\hyp Yapa languages), and in a prospective\slash future marker with and without negation (\citeauthor{chapters/fedotov} for \ili{Gban}; \citeauthor{chapters/song} for \ili{Baoding}). Pragmatic strengthening of an undesirability component is also involved in the development of general purpose/manner markers into precautionary markers, likewise discussed in this volume by \citeauthor{chapters/rhee-kuteva} for \ili{Korean} and by \citeauthor{chapters/treis} for \ili{Kambaata}.


An apprehensional strategy which perhaps most clearly conveys the undesirability to the speaker of an event is the use of an optative, i.e.\ a modal or construction which conveys the wish of the speaker, in combination with negation (\cite{Dobrushina2006}; see also \sectref{subsec:apprsubtypes}, \sectref{intro:functions}). In this volume, negative optatives  as the source of apprehensionals are discussed for \ili{Slavic} languages by \citeauthor{chapters/wiemer} and for \ili{Thulung} by \citeauthor{chapters/lahaussois}. In \ili{Thulung} the optative marker itself demonstrably originates in a lexical verb `be auspicious, be good'; Lahaussois also shows how a subordination of the negative optative construction can give rise to a precautionary construction.

\largerpage
Conditional constructions, including pseudo-conditional conjunction or disjunction, make a preemptive or enabling event explicit in relation to the event to be prevented, as in \REF{ex:intro24}, and thus have been discussed as an apprehensional strategy (e.g.\ \cites{auwera1986}{dominicyfranken2002}[35--40]{narrog2012}{clancyetal1997}{knoob2008}). \citet{chapters/lahaussois}  describes a conditional construction as one of the apprehensional strategies of \ili{Thulung}. 

\ea\label{ex:intro24}
\ea\label{ex:intro24a} \textbf{If} you do not start revising now, you might fail.
\ex\label{ex:intro24b} Start revising now \textbf{or else/otherwise} you're going to fail.
\ex\label{ex:intro24c} Keep on idling \textbf{and} you'll fail.
\z
\z

 A recurrent grammaticalisation source which gives rise to apprehensionals are ablative spatial cases or adpositions, represented in this volume in Upper Tanana \ili{Dene} \citep{chapters/lovick} and in some \ili{Oceanic} languages of Vanuatu \citep{chapters/francois}. In this case, the semantic change arises through metaphorical extension (spatial `away from' $>$ `avoidance').

Another apparently cross-linguistically recurrent source construction are temporal connectives like `before' (see \REF{ex:intro13b} in \sectref{subsec:othergramm}), or conversely `soon', `later', as illustrated in the \ili{Kriol} example in \REF{ex:intro1} and the \ili{German} examples in \REF{ex:intro17} and \REF{ex:intro18}. Equivalents are attested in \ili{Dutch} \citep{boogaart2009,boogaart2020}, in a number of traditional Australian languages, in a variety of \ili{Hawaiian} that was used for interethnic communication also known as Pidgin \ili{Hawaiian} (\citealt{roberts2013} cited in \citealt[283]{angeloschultzeberndt2016}) and in the Gyalrongic (\ili{Sino-Tibetan}) language \ili{Japhug} \citep[1419--1420]{jacques2021}. In this volume, \citet{chapters/francez} briefly discusses a modern \ili{Hebrew} marker with an origin in a temporal connective.

Other grammaticalisation sources for (components of) apprehensionals are more rarely discussed. Curiously, some languages make use of so-called alternative{\hyp}sensitive particles in their apprehensional constructions: \ili{German} \textit{nachher} `later/\textsc{appr}' tends to co-occur with the temporal additive/modal particle \textit{noch} `still' (see \REF{ex:intro1}), and in \ili{Cuzco Quechua}, warnings typically involve a contrastive marker as well as a possibility modal (fieldwork Faller). \citet{chapters/francez} shows how the temporal additive/modal particle \textit{od} `still' in \ili{Hebrew} on its own can invite apprehensional inferences, and invokes the same association of a subsequent event with an undesirable event that presumably facilitates the development from a temporal connective to an apprehensional marker.

Two further grammaticalisation paths not attested in the languages covered in the contributions to this volume will be mentioned here since they shed some light on how grammatical person restrictions may come into existence (see \sectref{subsec:othergramm}). The first case comes from \ili{Huallaga Huánuco Quechua}. \citet[304]{weber1989} shows how \textit{kuyra:} (from the original \ili{Spanish} attention verb \textit{(ten) cuidado} `be careful') plus a subordinate clause expressing an undesirable event gave way to two distinct constructions. One \REF{ex:intro26a} requires a same-subject adverbial subordinator (if the logical subject is 2\textsuperscript{nd} person), in line with the fact that \textit{kuyra:} was originally a main clause 2\textsuperscript{nd} person imperative to which the clause expressing the undesirable event is subordinated. The other \REF{ex:intro26b} requires a nominalised clause plus the comitative.


\ea \label{ex:intro26} 
\langinfo{Huallaga Huánuco Quechua}{\ili{Quechuan}; Peru}{\citealt[304]{weber1989}}\\
\ea \label{ex:intro26a} \textnormal{For 2\textsuperscript{nd} person subject} \\
\gll \textbf{Kuyra:} tuny-r. \\
\textbf{\textsc{appr}} fall-\textsc{\textbf{ss.adv}} \\
\glt `Be careful you might fall!' \\
\ex \label{ex:intro26b} \textnormal{For persons other than 2\textsuperscript{nd} person subject} \\
\gll \textbf{Kuyra:} pay	pusha-shu-\textbf{na}-yki-\textbf{wan}. \\
\textbf{\textsc{appr}} 3 lead-2-\textsc{\textbf{nmzr-}2p-\textbf{com}} \\
\glt `He may lead you astray.'
\z
\z

\noindent If one imagined a kind of Jespersen cycle where \textit{kuyra:} disappeared, one can see how two very distinct constructions may arise for different grammatical persons.

Another attested grammaticalisation path explaining a grammatical person restriction is that of the \ili{Matsés} apprehensive construction. \citet[443]{fleck2003} discusses how the Apprehensive \textit{-me-enda}, originally a Causative followed by a Prohibitive, is regularly shortened to \textit{-me}. Example \REF{ex:intro27} below is currently best translated as a proper apprehensive, but the full form would be literally translated as `don't let the jaguar kill you!'.


\ea \label{ex:intro27} 
\langinfo{Matsés}{\ili{Pano-Tacanan}; Brazil/Peru}{\citealt[443]{fleck2003}, our glosses}\\
\gll Bëdi ac-me. \\
jaguar kill-\textsc{caus/appr} \\
\glt `A jaguar might kill you.' 
\z

\noindent Interestingly, \citet[443]{fleck2003} describes this construction as being restricted to third person subjects, and most frequent with second person objects (first and third person objects are however grammatically possible). This restriction has a diachronic explanation, since a third person negatively affecting a second person is the most expected scenario for a warning, assuming that a warning use was the most frequent use of this construction at some point.

The bewildering variety of potential sources for apprehensionals just discussed may well be due to their particularly complex semantics (\cite{kutevaetal2019,anderboisdabkowskiupcoming}) and the interactional nature of their use. Possibly for the same reason, apprehensional constructions do not seem to be a diachronically stable phenomenon, an observation that is underlined by several contributions to this volume. The individual languages from the \ili{Ngumpin-Yapa} \citep{chapters/browne}, \ili{Oceanic} (\citealt{chapters/francois}; see also \cite{smithdennis2021}) and \ili{East Caucasian} \citep{chapters/daniel} language groups have quite distinct strategies for the expression of apprehensionality. The apprehensional markers in \ili{Gban} \citep{chapters/fedotov} and \ili{Kambaata} \citep{chapters/treis} appear to be recent innovations in these languages which are not shared in the wider family or area. Not uncommonly, languages have multiple co-existing strategies for conveying apprehensional meaning (\citealt{chapters/francez} for \ili{Hebrew}, \citealt{chapters/lahaussois} for \ili{Thulung}, \citealt{chapters/rhee-kuteva} for \ili{Korean},  \citealt{chapters/wiemer} for \ili{Slavic} languages). Apprehensional markers are also found in new languages i.e.\ creoles, e.g.\ \ili{Tok Pisin} (briefly discussed by \cite{smithdennis2021}) and Northern Australian \ili{Kriol} \citep{angeloschultzeberndt2016, phillips2021}.


We hope that this volume lays some groundwork for future research on the origins and the syntactic, semantic and pragmatic variety of apprehensional constructions in the world's languages. The last section concludes this introduction on the apprehensional domains covered in this volume with brief overviews of each of the contributions. They are geographically ordered from West to East, starting in North America and ending in Australia.

\section{Overview of this volume}
\label{sec:summaries}

In her paper \textit{Speaking and acting with care: The grammar of circumspection in \ili{Upper Tanana Dene},} \textsc{Olga} \textsc{Lovick} casts both a linguistic and anthropological perspective on two apprehensional constructions in \textsc{Upper} \textsc{Tanana} \textsc{Dene} (\ili{Athabaskan-Eyak-Tlingit}), spoken in the Alaska/Yukon border region between the US and Canada. The avertive clause construction in precautioning function involves the marker \textit{ch'à'} and the ``urgent request'' construction the marker \textit{dè'} (+ Optative). The origin of the avertive marker, which has cognates in the same function in several other \ili{Dene} languages, is likely the Proto-\ili{Dene} postposition *\textit{k'yaˑn} `away, apart from'. It exclusively encodes a direct causal relation with the preemptive main clause. Interestingly, the preemptive clause can never be negated. The ``urgent request'' construction is used to form expressions with translations like `make sure X happens', and likely originated in an insubordinated conditional request. The construction is linked to the apprehensional domain in that non-compliance with a request expressed in this way is likely to have severe  consequences, although they are never spelled out. Having discussed how traditional Upper Tanana speakers (and other Northern \ili{Dene} peoples) ``are very circumspect both in words and actions'', Lovick demonstrates how the two constructions allow the speaker to help the addressee prepare for or avoid undesirable eventualities while still respecting the strong cultural principle of non-interference.



\textsc{Maksymilian} \textsc{Dąbkowski}  and \textsc{Scott} \textsc{AnderBois} present a detailed unified analysis of the multi-functional apprehensional morpheme =\textit{sa'ne} in \textsc{\ili{A'ingae}} (aka.\ \ili{Cofán}), an isolate language spoken in the Northeastern Ecuador/Southern Colombia border region in Amazonia. Considering semantic and syntactic evidence as well as the paradigmatic status of this morpheme, they argue for a semantically conventionalised status of the undesirability component for =\textit{sa'ne.} Its primary function is that of a subordinator, i.e.\ a precautioning morpheme, with both avertive and `in-case' uses.  Dąbkowski and AnderBois propose a unified semantic analysis for these two uses of =\textit{sa'ne} (and comparable markers cross-linguistically): =\textit{sa'ne} always invokes an undesirable situation which is to be avoided. In the avertive reading this situation is the one described by the subordinate clause itself, whereas in the `in-case' reading it is a larger situation that contains the one described. The  marker moreover covers the three other apprehensional uses identified in the literature: it can function as a complementiser with `fear' verbs (and some other verbs with negative semantics), and has timitive uses with noun phrases which are interpretable as metonymically invoking an undesirable situation. It also has more marginal main-clause uses as an apprehensive, which are analysed as incipient insubordination.



\textsc{Maksim} \textsc{Fedotov} discusses the apprehensional construction in \textsc{\ili{Gban}}, a South \ili{Mande} language spoken in Côte d'Ivoire. This construction covers the three apprehensional core functions attested cross-linguistically. In independent clauses with an apprehensive function, the apprehensional construction is compatible with all three grammatical persons, in a pure apprehension, warning, or threat reading. In dependent clauses with a precautioning or `fear' complementation function, the construction requires additional subordinators. Fedotov argues that the origin of this formally complex construction is plausibly a negative purpose construction, in turn grammaticalised from the combination of a negator and a prospective morpheme, and that the precautioning construction still has  a residual negative semantics, as indicated by the availability of negative polarity items. The paper pays detailed attention to constraints on all three uses, e.g.\ that in the apprehensive function, the construction is always future-oriented and is never applied to events that cannot be in principle prevented, e.g.\ weather events (though it can be used for avoidable events involving weather events, such as rain falling on one's car).



\textsc{Yvonne} \textsc{Treis} discusses apprehensionality in \textsc{\ili{Kambaata}}, a \ili{Cushitic} language of Ethiopia, which seems to be the only \ili{Afroasiatic} language to display a verbal paradigm of dedicated apprehensive forms. These forms only occur in main clauses, and Treis provides detailed evidence showing that they do not exactly fit the verb forms of either the indicative aspectual paradigm or that of the directive. When the apprehensive clause co-occurs with a preemptive clause, the latter is either an independent clause or (unexpectedly from a cross-linguistic perspective) subordinated to the former. The apprehensive is available with all grammatical persons and shows cross-linguistically expected variation in the resulting pragmatic interpretation: a warning interpretation is available for all persons, while threat and prohibition uses are only attested for 1\textsuperscript{st} and 2\textsuperscript{nd} person, respectively. Treis methodically demonstrates that, like other verbal markers in the language, the apprehensive forms have a phrasal origin, which in this case involves a purposive converb and an existential copula.



\textsc{Itamar} \textsc{Francez{}'} overview of apprehension in \textsc{\ili{Hebrew}} covers five markers with apprehensional interpretations. Two of these (the complementiser \textit{pen} and the modal \textit{alul}) are argued to conventionally encode (or to have encoded at a previous diachronic stage) that the core proposition is possible and undesirable. Francez argues that the undesirability component is a conventional implicature, while the possibility meaning is an entailment. Both these markers can be used in the apprehensive function in matrix contexts, and in the precautioning function when the clause they appear in occurs in hypotaxis or parataxis with another clause. Apprehensional interpretations also arise with two temporal adverbials, one having an additive meaning (\textit{od} `still'), and the other a sequential meaning (\textit{axarkax} `afterwards/later'). The future orientation of both arguably provides the modal component of their apprehensional interpretations. The undesirability component however is not always present, and Francez therefore analyses it as a pragmatic inference that arises only in contexts which establish the associated clause as undesirable. Finally, a complex complementiser \textit{še lo –} consisting of a complementiser that can also stand on its own and a negation marker – has acquired an apprehensional interpretation, including when introducing the complement of `fear' predicates. Francez argues that negation in this construction is expletive, i.e.\ is no longer semantically interpreted.



In \textit{Two major apprehensional strategies in \ili{Slavic}: A survey of their areal and grammatical distribution}, \textsc{Björn} \textsc{Wiemer} offers a comprehensive survey of two apprehensional strategies widely attested in \textsc{Slavic} \textsc{languages}, and their genealogical and areal distribution in all three branches (East, West and South \ili{Slavic}). The discussion is based on data from both major corpora and linguistic publications often addressed to Slavic specialists. The first strategy foregrounds the speaker's apprehension of some undesirable event. It is dubbed the \textsc{negative-optative} strategy since it is mostly expressed by clause-initial connectives which are formally negated optatives, but it can also be realised by constructions associated with negative purpose clauses. The second strategy alerts the addressee to an action to be taken to avoid an undesirable situation. This apprehensional strategy includes what is known as \textit{preventives} in the literature on \ili{Slavic} languages but is more generally dubbed \textsc{negative-directive} in this paper. It most commonly involves a negated imperative (or equivalent clause type) with perfective verbs, whereas the imperative with imperfective verbs tends to express prohibition, at least in North \ili{Slavic}. Wiemer also offers an account of the diachronic development and grammaticalisation of the language-specific constructions and their synchronic interpretations.



\textsc{Michael} \textsc{Daniel} and \textsc{Nina} \textsc{Dobrushina}'s study considers dedicated apprehensionals in three distantly related \textsc{\ili{Nakh-Daghestanian}} (\ili{East Caucasian}) languages, \textsc{\ili{Archi}} (\ili{Lezgic}), \textsc{\ili{Mehweb}} (\ili{Dargwa}) and \textsc{Northern} \textsc{\ili{Akhvakh}} (\ili{Andic}). The markers in the three languages have distinct forms with different syntactic properties and probably arose through different pathways of grammaticalisation. The \ili{Akhvakh} morpheme exclusively expresses the apprehensive function, i.e.\ it is a ``discourse interactional category''. The \ili{Mehweb} morpheme primarily expresses the apprehensive function, i.e.\ it occurs in main clauses, but it extends to subordinate uses. In those uses, a (perfective) converb form of a `say' verb as subordinator is optional in precautioning function but obligatory in `fear' complement function. By contrast, the \ili{Archi} morpheme, for which the most data were available, primarily expresses the precautioning and `fear' complementation function, and its main clause functions seem to have developed from `fear' complement clauses via insubordination. The syntactic status of \ili{Archi} and \ili{Mehweb} apprehensional clauses is tested and discussed in detail, as well as their semantics, showing that at least in the usage of the \ili{Archi} apprehensional, the emotion of fear is not necessarily present, but only sheer undesirability. Whether the functional convergence of the apprehensional constructions in the three languages –~despite their different historical pathways – is the result of areal pressure remains to be investigated.



While \textsc{\ili{Thulung}} (\ili{Kiranti}, \ili{Sino-Tibetan}) has not developed a dedicated apprehensional construction – seemingly in line with other South Asian languages, which have been remarkably absent in the apprehensional literature so far – \textsc{Aimée} \textsc{Lahaussois}, in her paper \textit{Inauspicious events in \ili{Thulung},} explores how the negative purpose construction fulfills the precautioning function. The construction rather transparently originated in a negative optative verb form, which systematically associated with a subordinator typically introducing reported speech elsewhere in the language, originally a grammaticalised `say/think' converb. Such a developmental pathway is not only found in many Tibeto-Burman languages, but in the larger region. The optative itself goes back to a lexeme \textit{nʉmu} {}`be good, auspicious', which gave rise to two distinct grammatical constructions, the impersonal negative deontic necessity construction on the one hand, and the precautionary construction on the other. This article provides useful information on how a language with no dedicated apprehensional construction may develop a strategy to express the function. It also convincingly links the semantics of the original lexical concept with apprehension and deontic necessity, and highlights its conceptual relation with \ili{Thulung} cultural practices.


\textsc{Na} \textsc{Song} examines the apprehensive function of a sentence-final particle in \textsc{\ili{Baoding}}, a \ili{Mandarin} (\ili{Sinitic}) dialect spoken in Northern China. The Apprehensive enclitic \textit{=lɛ.ia} is atypical for an apprehensive in several regards: syntactically, it can occur in interrogative clauses; semantically, it largely conveys unexpectedness (on top of undesirability); and pragmatically, it requires the addressee's involvement in avoiding the potential undesirable situation. The particle systematically expresses apprehension with change of state verbs, and tends to express imminent future with activity verbs. The grammatical first person however seems to force an apprehensive reading in the latter case. Song provides phonological, morphosyntactic, semantic and typological arguments for a diachronic source of the marker in the univerbation of two final particles, =\textit{lɛ} denoting a change of state, and =\textit{ia} denoting the intention of the speaker about an imminent future action.

\textsc{Songha} \textsc{Rhee} and \textsc{Tania} \textsc{Kuteva's} contribution addresses the diachronic development and synchronic distribution of no fewer than 13 markers in apprehensional function in \textsc{Korean} (also known as preventives in \ili{Korean} linguistics). Using both historical and contemporary corpora, Rhee and Kuteva are able to trace the first attestations of these markers as well as differences in their frequency and in their use over time, including the degree to which each marker appears to be restricted to apprehensional contexts. They therefore posit multiple layers in the functional domain of apprehensionality for present-day \ili{Korean}. They also provide an in-depth discussion of the grammaticalisation processes that gave rise to each marker, and argue that some of the markers preserve a degree of semantic transparency and analysability (lack of erosion) due to not being fully grammaticalised. The source constructions of one set of markers – in a predominantly precautionary function – involve combinations of a nominaliser and either a mode/manner or purposive marker. The source of the others are various combinations of lexical verbs expressing fear but also perception (`see') and cognition (`not know', `think') with existing grammaticalised future and interrogative markers. The grammaticalisation of perception and cognition verbs in these contexts provides strong support for a cognitive link between uncertainty about the future and anxiety, which can lead to the development of apprehensional markers from expressions of mere possibility via pragmatic strengthening. These forms are also often ambiguous between a main clause use and a subordinate/connective use, which is attributed to processes of insubordination and incoordination, also attested elsewhere in the history of \ili{Korean}. Interestingly, for those markers, the apprehensional function is more prominent in the corpus attestations in their subordinate (connective) use than in their main clause use.

In \textit{Explicit apprehensions, implicit instructions: An indirect speech act in the grammar,} \textsc{Alexandre} \textsc{François} offers a survey of apprehensional strategies in the \textsc{\ili{Oceanic}} \textsc{languages} of north Vanuatu, followed by an in-depth study of the syntactic, semantic and pragmatic properties of a dedicated apprehensive marker in one of these languages, \textsc{\ili{Mwotlap}}. Despite their close genealogical relationship, the languages of the Torres-Banks area of north Vanuatu exhibit a fascinating diversity of apprehensionals (subordinators and/or modal markers), with interesting patterns of polyfunctionality in that they overlap with either an ablative case marker or a prohibitive marker (and in the case of \ili{Lemerig}, with both). Out of 8 languages surveyed, only \ili{Mwotlap} has a distinct apprehensive mood marker (\textit{tiple} and allomorphs). Clauses featuring \textit{tiple} can function as complements of `fear' verbs, and are also frequently found in combination with a preemptive clause in both avertive and `in-case' functions. However, François argues that \textit{tiple} is not a subordinator and that the relationship of apprehensive clauses with preemptive clauses is one of pragmatic rather than syntactic dependency. Thus, apprehensive clauses on their own can be exploited to convey an implicit instruction (that of taking appropriate action to avoid the situation described by the apprehensive clause), i.e.\ an indirect speech act, which just like other cases of indirect speech act serves as a politeness strategy: an explicit directive can be avoided by merely pointing out to the addressee the consequences of not taking action.


\textsc{Mitchell} \textsc{Browne}, \textsc{Thomas} \textsc{Ennever} and \textsc{David} \textsc{Osgarby} explore both dedicated and non-dedicated apprehensional expressions at all levels of grammar in a dozen languages of the \textsc{\ili{Ngumpin-Yapa} subgroup} of \ili{Pama-Nyungan} (Australia), in a comparative Australianist perspective. Despite their relatively close genealogical relationship, the languages examined exhibit an interesting diversity in the apprehensional domain. Most \ili{Ngumpin-Yapa} languages surveyed have an aversive case suffix with the specialised function of indicating that a nominal referent is to be feared or avoided. This can have a subordinating function in non-finite clauses, and thus form precautioning subordinate constructions. It can also mark (nominal or non-finite clause) complements of `fear' predicates in some languages, and even displays insubordinate uses with the illocutionary force of a directive. The second major strategy for encoding apprehensional meanings in \ili{Ngumpin-Yapa} languages involves grammatical markers variously analysed as auxiliaries or complementisers/subordinators in existing descriptions. The authors favour an analysis as auxiliaries, even when a preemptive clause is mostly present in the context, i.e.\ they analyse the combination of preemptive clause and apprehensional clause as cases of juxtaposition, with a pragmatic rather than syntactic relationship between the two, although they consider the potential of a diachronic development to subordination constructions. The semantic interpretation of the auxiliaries may also depend on the verbal inflection, e.g.\  apprehension with future inflection but past counterfactual with past inflection. The authors also demonstrate that such auxiliaries are dedicated apprehensives in some of the languages surveyed but have a broader irrealis, habitual or capability meaning in others. Such semantic specialisation likely arose through pragmatic strengthening, from a general possibility modal to an apprehensive marker.


\textsc{Marine Vuillermet}'s contribution differs significantly from the others as it is a questionnaire intended as a guide to fieldworkers. It systematically explores the semantic, syntactic and pragmatic specificities of the four main construction types in the apprehensional domain – namely apprehensive, precautioning, `fear' complementation and timitive. It outlines the potential variation in the properties of the respective constructions that is expected based on a cross-linguistic survey of apprehensionals, and provides contexts targeting the corresponding parameters in the form of over 80 prompts in \ili{English}, to be translated into the language of study (and/or adapted to the cultural context). It is hoped that this questionnaire will enable the systematic investigation of the fine-grained semantics of apprehensional constructions, which will in turn contribute to a better understanding of the intricacies of the apprehensional functional domain.

\section*{Acknowledgements}
We are grateful to Ma\"{\i}a Ponsonnet for her valuable comments on a draft version of this chapter. We would also like to give huge thanks to Siena Weingartz for her thorough attention to detail in converting chapters to \LaTeX\ and copy-editing. We also thank Matthew Korte for his thorough proofreading of the entire volume, and Sebastian Nordhoff and Felix Kopecky at Language Science Press for their help with typesetting the final version. Part of the data for this paper have been gathered during Vuillermet's postdoctoral fellowship from LABEX ASLAN (ANR-10-LABX-0081) of the Université de Lyon within the program ``Investissements d'Avenir'' (ANR-11-IDEX-0007).


\section*{Abbreviations}
\begin{tabularx}{.5\textwidth}{@{}lQ@{}}
1/2/3 & first/second/third person \\
\textsc{a} & agent \\
\textsc{abs} & absolutive \\
\textsc{add} & additive \\
\textsc{adv} & adverb \\
\textsc{appr} & apprehensive \\
\textsc{caus} & causative \\
\textsc{cntr} & control \\
\textsc{com} & comitative \\
\textsc{dat} & dative \\
\textsc{decl} & declarative \\
\textsc{def} & definite \\
\textsc{erg} & ergative \\
\textsc{expl.neg} & expletive negation \\
\textsc{fact} & factive \\
\textsc{fem} & feminine \\
\textsc{fut} & future \\
\textsc{imp} & imperative \\
\textsc{incl} & inclusive \\
\textsc{ind} & indicative \\
\textsc{inst} & instrumental \\\\
\end{tabularx}%
\begin{tabularx}{.5\textwidth}{@{}lQ@{}}
\textsc{irr} & irrealis \\
\textsc{mid} & middle \\
\textsc{neg} & negation \\
\textsc{nf} & non-feminine \\
\textsc{nfut} & non-future \\
\textsc{nmzr} & nominaliser \\
\textsc{obj} & object \\
\textsc{p} & patient \\
\textsc{pl} & plural \\
\textsc{prec} & precautioning \\
\textsc{proh} & prohibitive \\
\textsc{propr} & proprietive (`having') \\
\textsc{prs} & present \\
\textsc{sbjv} & subjunctive \\
\textsc{seq} & sequential \\
\textsc{sg} & singular \\
\textsc{ss} & same subject \\
\textsc{sub} & subordinator  \\
\textsc{tim} & timitive \\
\textsc{tmp} & temporal subordinate \\
\textsc{vis} & visual evidential \\
\\
\end{tabularx}


\sloppy
\printbibliography[heading=subbibliography,notkeyword=this]
\end{document}
